\documentclass[12pt,a4paper]{amsart}
\usepackage{fontenc,amsmath,mathrsfs, amsthm, verbatim,amssymb, MnSymbol, enumerate, multicol, hyperref, geroENG, color}
\usepackage[english]{babel}
\allowdisplaybreaks
\usepackage{times}
%\usepackage{tgbonum}
\allowdisplaybreaks
\usepackage[top=2.25cm, bottom=2.25cm, left=2.5cm, right=2.5cm]{geometry}
\pagestyle{headings}
\usepackage{parskip}
\setlength{\parskip}{2ex}
\begin{document}
\title[Measure theoretic properties of large products]{Measure theoretic properties of large products of consecutive partial quotients}
\author[Brown-Sarre]{Adam Brown-Sarre}
\author[Gonz\'alez Robert]{Gerardo Gonz\'alez Robert}
\author[Hussain]{Mumtaz Hussain}
\address{Department of Mathematical and Physical Sciences,  La Trobe University, Bendigo 3552, Australia. }
\email[ABS]{20356213@students.latrobe.edu.au}
\email[GGR]{g.robert@latrobe.edu.au, matematicas.gero@gmail.com}
\email[MH]{m.hussain@latrobe.edu.au}
\subjclass[2020]{Primary 11K55; Secondary 11J83, 28A80}

\begin{abstract}
The theory of uniform approximation of real numbers motivates the study of products of consecutive partial quotients in regular continued fractions. For any non-decreasing positive function $\varphi:\mathbb{N}\to [2,\infty)$, we determine the Lebesgue measure and Hausdorff dimension of the set $\mathcal{F}_{\ell}(\vphi)$ of irrational numbers $x$ whose regular continued fraction $x~=~[a_1(x),a_2(x),\ldots]$ is such that for infinitely many $n\in\Na$ there are two numbers $1\leq j<k \leq n$ satisfying
\[
a_{k}(x)a_{k+1}(x)a_{k+2}(x) \geq \varphi(n), \;
a_{j}(x)a_{j+1}(x)a_{j+2}(x) \geq \varphi(n). 
\]
One of the consequences of the results is that the strong law of large numbers for products of $3$ consecutive partial quotients is impossible even if the block with the largest product is removed.
\end{abstract}
\maketitle

\section{Introduction}


Regular continued fractions are a representation system of real numbers with outstanding approximation properties. Recall that, for each $x\in \RE\setminus\QU$ there is a single sequence of integers $(a_j(x))_{j\geq 0}$ with $a_j(x)\geq 1$ for $j\geq 1$ such that
\[
x
=
a_0(x) + \cfrac{1}{a_1(x) + \cfrac{1}{a_2(x) + \cfrac{1}{\ddots}}}:=[a_0(x); a_1(x),a_2(x),\ldots].
\]
The numbers $a_n(x)$ are known as the partial quotients of $x$. We have $x\in [0,1)$ if and only if $a_0(x)=0$. In this case, we write the continued fractions of $x$ as $[a_1(x),a_2(x),\ldots]$. It is well known that the $n$-th convergent of $x$, given by
\[
\frac{p_n(x)}{q_n(x)}
=
[a_1(x),\ldots, a_n(x)]
\quad\text{ for all } n\in\Na,
\]
where $q_n(x)\geq 1$ and $p_n(x)$ are coprime integers, satisfies 
\[
\frac{1}{(a_{n+1}(x)+2)q_n^2(x)}
<
\left| x - \frac{p_n(x)}{q_n(x)}\right|
<
\frac{1}{q_n^2(x)}
\quad\text{ for all } n\in\Na,
\]
(see, for example, \cite[Equation (34)]{KhinchinCFBook}). Because  $(q_n(x))_{n\geq 1}$ is strictly increasing, the rate at which rationals approximate $x$ is closely related to the size of its partial quotients. This leads us to consider, for any function $\vphi:\Na\to\RE_{>0}$, the set
\[
\clE_1(\varphi)
:= \{x \in [0,1):a_n(x) \geq \varphi(n) \text{ for i.m. } n\in \mathbb{N}\}.
\]
Here and in what follows, ``i.m.'' stands for ``infinitely many''. The well-known Borel-Bernstein Theorem establishes the Lebesgue measure $\leb$ of $\clE_1(\vphi)$:
\[
\leb\left( \clE_1(\vphi) \right)
=
\begin{cases}
    0, \text{ if }  \displaystyle\sum_{n\geq 1} \frac{1}{\vphi(n)}<\infty,\\
    1, \text{ if }  \displaystyle\sum_{n\geq 1} \frac{1}{\vphi(n)}=\infty.\\
\end{cases}
\]
For proof of this theorem, see, for example, \cite[Theorem 1.11]{BugeaudBookCUP160}. {\L}uczak \cite{Luc1997} and Feng et al. \cite{FengWuLiang1997} calculated the Hausdorff dimension, denoted as $\hdim$, of $\clE_1(\vphi)$ when $\vphi(n)=c^{b^n}$ for some fixed $c,b>1$ (see Lemma \ref{Lem:Luczak}). Wang and Wu solved the general case in \cite{WangWu2008} (see Theorem \ref{TEO:HWX:HD} below with $\ell=1$).

While studying improvements to the classical Dirichlet theorem, Kleinbock and Wadleigh \cite{KleinbockWadleigh2019} showed that the set of Dirichlet non-improvable numbers is related to the set
\[
\clE_{2}(\varphi)
:= \{x \in [0,1):a_n(x)a_{n+1}(x) \geq \varphi(n) \text{ for i.m.   } n\in \mathbb{N}\}.
\]
In fact, they computed the Lebesgue measure of $\clE_{2}(\varphi)$ and the Hausdorff measure and dimension were proven in \cite{BHS, HKWW}. In \cite{HuangWuXu2020}, Huang, Wu, and Xu determined the Lebesgue measure and the Hausdorff dimension of the general set $\clE_{\ell}(\varphi)$, $\ell\in\Na$, defined as:
\[
\clE_{\ell}(\varphi)
:= \{x \in [0,1):a_n(x)\cdots a_{n+\ell-1}(x) \geq \varphi(n) \text{ for i.m.   } n\in \mathbb{N}\}.
\]
\begin{teo01}[Huang, Wu, Xu \cite{HuangWuXu2020}]\label{TEO:HWX:MEAS}
If $\vphi:\Na\to \RE_{\geq 2}$, then
\[
\leb\left( \clE_{\ell}(\vphi)\right)
=
\begin{cases}
0, &\text{\rm if } \displaystyle\sum_{n=1}^{\infty} \frac{\log^{\ell-1}\vphi(n)}{\vphi(n)}<\infty, \\    
1, &\text{\rm if } \displaystyle\sum_{n=1}^{\infty} \frac{\log^{\ell-1}\vphi(n)}{\vphi(n)}=\infty.
\end{cases}
\]
\end{teo01}

Define $B$ and $b$ as
\begin{equation}\label{Bb}
\log B := \liminf_{n \xrightarrow{} \infty} \frac{\log \varphi(n)}{n}, 
\;\text{ and }\;
\log b = \liminf_{n \xrightarrow{} \infty} \frac{\log \log \varphi(n)}{n}.
\end{equation}
\begin{teo01}[Huang, Wu, Xu \cite{HuangWuXu2020}]\label{TEO:HWX:HD}
Let $(f_n)_{n\geq 1}$ be the sequence of real functions on $[0,1]$ given by
\[
f_1(s)=s, 
\;\text{ and }\;
f_{n+1}=\frac{sf_n(s)}{1-s+f_n(s)} \;\text{ for all } n\in\Na.
\]
Let $\vphi:\Na\to\RE_{>0}$ be arbitrary and $B$, $b$ as in \eqref{Bb}. Then, the Hausdorff dimension of $\clE_{\ell}(\vphi)$ is 
\[
\hdim \clE_{\ell}(\vphi)
=
\begin{cases}
1, & \text{\rm if } B=1, \\
\inf\{s \geq 0: P(T,-f_{\ell}\log B - s \log |T'|) \leq 0\},  & \text{\rm if }  1<B<\infty,\\
\frac{1}{1+b},  & \text{\rm if }  B=\infty.
\end{cases}
\]
\end{teo01}
The pressure function $P(T,\cdot)$ appearing in the Hausdorff dimension formula above and in all the theorems below will be defined in Section \ref{Section:Preleminaries}.


From a probabilistic perspective, we may interpret partial quotients as random variables defined on $[0,1)$ (their definition of rational numbers is irrelevant for our purposes). These random variables are not independent, hence the Borel-Bernstein Theorem is not a trivial consequence of the Borel-Cantelli lemma. Moreover, the Borel-Bernstein theorem tells us that the inequality $a_n(x)\geq n\log n$ holds infinitely often for almost every number $x$ (with respect to $\leb$). This implies that, writing $S_n(x):=a_1(x) + \cdots + a_n(x)$, we have
\[
\limsup_{n\to\infty} \frac{1}{n} S_n(x) = \infty
\;\text{ for almost all } x\in [0,1).
\]
Despite not being independent and having infinite expected value, the partial quotients satisfy a weak law of large numbers. Khinchin \cite[\S. 4]{Khintchine1935} proved that $S_n/(n\log n)$ converges in Lebesgue measure to $\frac{1}{\log 2}$ as $n\to\infty$. Furthermore, Philipp \cite[Theorem 1]{Philipp1988} proved that there is no reasonable function $\sigma$ for which $S_n/\sigma(n)$ converges almost everywhere to a finite positive constant. Large partial quotients are to blame for this phenomenon. Certainly, Diamond and Vaaler \cite[Corollary 1]{DiamondVaaler1986} showed that convergence almost everywhere does hold for Khinchin's averaging function when omitting the largest partial quotient: for almost every $x\in [0,1)$ we have
\begin{equation}\label{Eq:Thm:DV}
\lim_{n\to\infty} 
\frac{1}{n\log n}\left( S_n(x) - \max_{1\leq j \leq n} a_j(x)\right)
=
\frac{1}{\log 2}.
\end{equation}
This result was extended by Hu, Hussain, and Yu for sums of products of two partial quotients \cite[Theorem 1.5]{HuHussainYu2021}. They showed that almost every $x\in [0,1)$ satisfies
\begin{equation}\label{Eq:Thm:HHY}
\lim_{n\to\infty} 
\frac{1}{n\log^2 n}\left( \sum_{j=1}^n a_j(x)a_{j+1}(x) - \max_{1\leq j \leq n} a_j(x)a_{j+1}(x)\right)
=
\frac{1}{2\log 2}.
\end{equation}

A key aspect in the proofs by Diamond and Vaaler, and Hu, Hussain, and Yu respectively is proving that, for $c>\frac{1}{2}$ and $d>\frac{3}{2}$, the sets
\[
\left\{
x\in [0,1):
\begin{matrix}
a_j(x) \geq  n (\log n)^{c}, \\
a_k(x) \geq  n (\log n)^{c},
\end{matrix}\;
\text{ for some } 1\leq j < k \leq n \text{ for i.m.   } n\in\Na
\right\}
\]
and
\[
\left\{x\in [0,1): 
\begin{matrix}
a_j(x)a_{j+1}(x) \geq n (\log n)^{d} , \\
a_k(x)a_{k+1}(x) \geq n (\log n)^{d},  \\
\end{matrix}\;
\text{ for some } 1\leq j < k \leq n \text{ for i.m.   } n\in\Na\right\}
\]
are of Lebesgue measure $0$.

Recently, Tan, Tian, and Wang \cite{TanTianWang2023} generalized this result of Diamond and Vaaler to a wider family of functions. Later, Tan and Zhou \cite{TanZhou2024} also extended the result of Hu, Hussain and Yu. More precisely, define 
\[
\clF_{\ell}(\vphi)
:=
\left\{ x\in [0,1):
\begin{matrix}
a_{j}(x)\cdots a_{j+\ell-1}(x) \geq \vphi(n), \\
a_{k}(x)\cdots a_{k+\ell-1}(x) \geq \vphi(n),
\end{matrix}\;
\text{ for some } 1\leq j < k \leq n \text{ for i.m.   } n\in\Na\right\}.
\]
In both \cite{TanTianWang2023} and \cite{TanZhou2024}, $\vphi$ is assumed to be non-decreasing. It is obvious that $\clF_{\ell}(\vphi)=[0,1)\setminus \QU$ whenever $\vphi(n)\leq 1$ for all $n\in \Na$. Hence, in what follows we suppose that $\vphi$ is non-decreasing and $\vphi:\Na\to[2,\infty)$.
 
\begin{teo01}[Tan, Tian, Wang \cite{TanTianWang2023}]\label{TEO:TTW:MEAS} Let  $\vphi:\Na\to[2,\infty)$ be non-decreasing. Then, the Lebesgue measure of $\clF_1(\vphi)$ is 
    \[
    \leb(\clF_1(\vphi))
    =
    \begin{cases}
        0, &\text{\rm if }  \displaystyle\sum_{k\geq 1} \frac{n}{\varphi(n)^2}<\infty,\\
        1, &\text{\rm if }  \displaystyle\sum_{k\geq 1} \frac{n}{\varphi(n)^2}=\infty.
    \end{cases}
    \]
\end{teo01}

%In \cite{TanZhou2024}, Tan and Zhou computed the Lebesgue measure of $\clF_2(\vphi)$.
\begin{teo01}[Tan, Zhou, 2024]\label{TEO:TZ:MEAS}
%Let $\vphi:\Na\to\Na$ be a non-decreasing function such that $\vphi(n)\to\infty$ as $n\to\infty$. 
 Let  $\vphi:\Na\to[2,\infty)$ be non-decreasing. Then, the Lebesgue measure of $ \clF_2(\vphi)$ is given by
\[
\leb(\clF_2(\vphi))
=
\begin{cases}
0, &\text{\rm if } \displaystyle\sum_{n\geq 1} 
\frac{n\log  \vphi(n)}{\vphi^2(n)} + \frac{1}{\vphi(n)} <\infty, \\   1, &\text{\rm if } \displaystyle\sum_{n\geq 1} 
\frac{n\log  \vphi(n)}{\vphi^2(n)} + \frac{1}{\vphi(n)} =\infty.
\end{cases}
\]
\end{teo01}
%Our first result, which includes Theorem \ref{TEO:TZ:MEAS}, completes the picture of the Lebesgue measure of $\clF_{\ell}(\vphi)$. 
Our first result extends Theorem \ref{TEO:TZ:MEAS} to blocks of length $3$.

%\begin{teo01}\label{TEO:MAIN:MEASURE}
% Let  $\vphi:\Na\to[2,\infty)$ be non-decreasing. For $\ell\geq 2$, the Lebesgue measure of $ \clF_{\ell}(\vphi)$ is given by
%\[
%\leb(\clF_{\ell}(\vphi))
%=
%\begin{cases}
%0 &\text{\rm if } \displaystyle\sum_{n\geq 1} 
%\frac{n\log^{2(\ell-1)}\vphi(n)}{\vphi^2(n)} + \frac{\log^{2(\ell-2)}\vphi(n)}{\vphi(n)}<\infty, \\
%1 &\text{\rm if } \displaystyle\sum_{n\geq 1} 
%\frac{n\log^{2(\ell-1)}\vphi(n)}{\vphi^2(n)} + \frac{\log^{\ell-2}\vphi(n)}{\vphi(n)}=\infty.
%\end{cases}
%\]
%\end{teo01}
%{\color{red}
\begin{teo01}\label{TEO:MAIN:MEASURE}
%Let $\vphi:\Na\to\Na$ be a non-decreasing function. 
 Let  $\vphi:\Na\to[2,\infty)$ be non-decreasing. 
 The Lebesgue measure of $ \clF_{3}(\vphi)$ is given by
\[
\leb(\clF_{3}(\vphi))
=
\begin{cases}
0 &\text{\rm if } \displaystyle\sum_{n\geq 1} 
\frac{n\log^{4}\vphi(n)}{\vphi^2(n)} + \frac{\log \vphi(n)}{\vphi(n)}<\infty, \\
1 &\text{\rm if } \displaystyle\sum_{n\geq 1} 
\frac{n\log^{4}\vphi(n)}{\vphi^2(n)} + \frac{\log \vphi(n)}{\vphi(n)}=\infty.
\end{cases}
\]
\end{teo01}
%}
\begin{rema01}
Using Birkhoff's Ergodic Theorem, we may show that for almost every $x\in [0,1)$ we have
\[
\lim_{n\to\infty}
\frac{1}{n}
\sum_{j=1}^n a_j(x)\cdots a_{j+\ell}(x)
=
\infty
\]
(see, for example, \cite[Corollary 3.8]{EinsiedlerWardET}).
Moreover, Theorem \ref{TEO:MAIN:MEASURE} implies that for almost every $x\in [0,1)$ there are infinitely many numbers $n\in\Na$ for which there are $1\leq j<k\leq n$ satisfying 
\[
a_{j}(x)a_{j+1}(x)a_{j+2}(x)  \geq n\log^{2}n
\;\text{ and }\;
a_{k}(x)a_{k+1}(x)a_{k+2}(x)  \geq n\log^{2}n.
\]
Therefore, we have 
\[
\limsup_{n\to\infty}  
\frac{1}{n} \left(\sum_{j=1}^n a_{j}(x)a_{j+1}(x)a_{j+2}(x)  
- 
\max_{1\leq k\leq n} a_{k}(x)a_{k+1}(x)a_{k+2}(x) \right)
=
\infty.
\]
Theorem \ref{TEO:MAIN:MEASURE} is also a crucial step in an extension of the limits \eqref{Eq:Thm:DV} and \eqref{Eq:Thm:HHY}. In fact, in a forthcoming paper we rely on Theorem \ref{TEO:MAIN:MEASURE} to show that for almost every $x\in [0,1)$ we have 
\[
\lim_{n\to\infty} 
\frac{1}{n\log^{3} n} 
\left( \sum_{j=1}^na_{j}(x) a_{j+1}(x)a_{j+2}(x) 
- 
\max_{1\leq k\leq n} a_{k}(x)a_{k+1}(x)a_{k+2}(x)\right)
=
\frac{1}{3 \log 2}.
\]
Other anticipated results are stated in Section \ref{Section:FinalRemarks}. 

\end{rema01}

%{\color{red} As discussed with you, we should add a full section on claiming the proofs of limit theorems for arbitrary product of consecutive and partial quotients from arithmetic progressions, maybe in the last section. The limit theorems will justify the existence of max terms in the large partial quotients}

\begin{rema01}
While proving Theorem \ref{TEO:MAIN:MEASURE} we must obtain new estimates corresponding to overlapping blocks of consecutive partial quotients with large products. The number of terms overlapping in Theorem \ref{TEO:MAIN:MEASURE} ranges from $0$ up to $\ell-1=2$ whereas in Theorem \ref{TEO:TZ:MEAS} the only overlaps are of length $0$ and $1$ and there are no overlaps in Theorem \ref{TEO:TTW:MEAS}.
\end{rema01}

The Hausdorff dimensions of $\clF_{1}(\vphi)$ and $\clF_{2}(\vphi)$ were computed, respectively, in \cite{TanTianWang2023} and \cite{TanZhou2024}. %In the three theorems to follow, $B$ and $b$ are defined as in Theorem \ref{TEO:HWX:HD}.

\begin{teo01}[Tan, Tian, Wang \cite{TanTianWang2023}]\label{TEO:TTW:HD}
Let $\vphi:\Na\to [2,\infty)$ be non-decreasing and $B$ and $b$ as in \eqref{Bb}. Then, the Hausdorff dimension of $\clF_1(\vphi)$ is given by
\[
\hdim \clF_1(\vphi)  
=
\begin{cases}
1, &\text{\rm if } B=1, \\
\inf\{s \geq 0: P(T,-(3s-1)\log B - s \log |T'|) \leq 0\}, &\text{\rm if }  1<B<\infty,\\
\frac{1}{1+b}, &\text{\rm if }  B=\infty.
\end{cases}  
\]
\end{teo01}

\begin{teo01}[Tan, Zhou \cite{TanZhou2024}]\label{TEO:TZ:HD}
%Let $\vphi:\Na\to\RE_{>0}$ be a non-decreasing function such that $\vphi(n)\to\infty$ when $n\to \infty$. 
Let $\vphi:\Na\to [2,\infty)$ be non-decreasing and $B$ and $b$ as in \eqref{Bb}.
%Let $\vphi$ be defined as in \eqref{Bb}.  
Then the Hausdorff dimension of  $\clF_2(\vphi)$ is given by
\[
\hdim \clF_2(\vphi)  
=
\begin{cases}
1, &\text{\rm if }  B=1, \\
\inf\{s \geq 0: P(T,-(3s-1-s^2)\log B - s \log |T'|) \leq 0\}, &\text{\rm if }  1<B<\infty,\\
\frac{1}{1+b}, &\text{\rm if } B=\infty.
\end{cases}  
\]
\end{teo01}
The case $1<B<\infty$ is the main substance of theorems \ref{TEO:TTW:HD} and \ref{TEO:TZ:HD}. In both theorems, this case can be obtained as an application of a general criterion established in \cite{ShulgaHussain2023Exp} (see Lemma \ref{Le:MumtazNikita}). We refer to  \cite{ShulgaHussainPisa} for the study of the product of partial quotients arising from arithmetic progressions and the metrical study of this set providing information on the size of the set of exceptions to the convergence criteria of multiple ergodic averages for the Gauss dynamical systems. 


Our second result provides the Hausdorff dimension of $\clF_3(\vphi)$. Let $g_3:\RE\to\RE$ be given by
\[
g_3(s)
:=
\frac{3s^3-5s^2+4s-1}{s^2-s+1}.
\]

\begin{teo01}\label{TEO:MAIN:HD}
Let $\vphi:\Na\to [2,\infty)$ be non-decreasing and $B$, $b$ as in \eqref{Bb}. Then, the Hausdorff dimension of $\clF_3(\vphi)$ is given by
\[
\hdim \clF_3(\vphi)  
=
\begin{cases}
1, &\text{\rm if }  B=1, \\
\inf\{s \geq 0: P(T,-g_3(s)\log B - s \log |T'|) \leq 0\}, &\text{\rm if }  1<B<\infty,\\
\frac{1}{1+b}, &\text{\rm if }  B=\infty.
\end{cases}  
\]
\end{teo01}



The paper is organized as follows. Section \ref{Section:Preleminaries} contains some preliminaries and auxiliary lemmas from measure theory and continued fractions, in particular on the Hausdorff dimension and pressure functions. In Section \ref{Section:Series} we estimate certain series that are fundamental for our proofs. The proof of Theorem \ref{TEO:MAIN:MEASURE} is in Section \ref{Section:ProofMeasure} and the proof of Theorem \ref{TEO:MAIN:HD} is in Section \ref{Section:ProofHD}. Lastly, Section \ref{Section:FinalRemarks} contains some final remarks and results developed to further this study.  

\paragraph{Notation.} Whenever $(x_n)_{n\geq 1}$ and $(y_n)_{n\geq 1}$ are two sequences of positive real numbers, we write $x_n\ll y_n$ to mean that for some constant $c>0$ and every $n\in\Na$ we have $x_n\leq cy_n$. When $c$ depends on some parameter $\lambda$, we write $x_n\ll_{\lambda} y_n$. By $x_n\asymp y_n$ we mean $x_n\ll y_n$ and $y_n\ll x_n$. For the rest of the paper, $\ell$ will denote an arbitrary natural number.


\medskip

\noindent{\bf Acknowledgements.} 
GGR and MH were funded by the Australian Research Council discovery project grant number 200100994. We thank Nikita Shulga, Valeriia Starichkova, and  Ben Ward for useful discussions.
\section{Preliminaries}\label{Section:Preleminaries}
In this section, we recall some definitions and elementary results on measure theory, Hausdorff dimension, continued fractions, and pressure functions. 



\subsection{Full and null sets}

The Borel-Cantelli Lemma is a fundamental tool for determining the measure of limsup sets. A proof may be found, for example, in \cite[Lemma 4.18]{KallenbergBook2021}. 

\begin{teo01}[Borel-Cantelli Lemma]\label{Le:BorelCantelli}
Let $(E_k)_{k\geq 1}$ be a sequence of measurable sets in a probability space $(\Omega,\mathscr{A},\mu)$ and $E=\limsup_n E_n$. 
\begin{enumerate}[\rm (i).]
\item If $\sum_{n}\mu(E_n)<\infty$, then $\mu(E)=0$.
\item If the collection $\{E_n:n\in\Na\}$ is stochastically independent and $\sum_{n}\mu(E_n)=\infty$, then $\mu(E)=1$.
\end{enumerate}
\end{teo01}
The independence assumption is rather restrictive, it will not hold in our context. A way to circumvent this problem is proving the positive measure of the limsup set and showing by some other means that the measure is either full or null. We show the full measure statement in Theorem \ref{TEO:MAIN:MEASURE} using a ubiquity-type property and Lemma \ref{Le:Knopp}. This lemma is attributed to Knopp \cite{Knopp1926} by Tan and Zhou \cite{TanZhou2024}. 

Given two measurable sets $A,B$ in a probability space $(X,\scA,\mu)$, we write $A \equiv B \pmod \mu$ if and only if $\mu(A\setminus B)=\mu(B\setminus A)=0$.

\begin{lem01}\label{Le:Knopp}
Let $A \subseteq [0,1)$ be a Borel set and $\mathcal{C}$ a class of subintervals of $[0, 1)$ satisfying both of the following conditions:
    \begin{enumerate}
        \item Every open subinterval of $[0,1)$ is an at most countable union $\pmod \leb$ of disjoint elements from $\mathcal{C}$,
        \item There is a constant $c>0$ such that for any $B \in \mathcal{C}$ we have $\leb (A\cap B) \geq c \leb(B)$.
    \end{enumerate}
    Then, $\leb(A)=1$.
\end{lem01}

In order to apply Lemma \ref{Le:Knopp}, we need the following result due to Chung and Erd{\H o}s \cite{ChungErdos1952} (see also \cite[Lemma 6]{Sprindzuk1979} for the formulation stated below).

\begin{lem01}[Chung, Erd{\H o}s, \cite{ChungErdos1952}]\label{Le:ErdosChung}
Let $(E_k)_{k\geq 1}$ be a sequence of measurable sets in a probability space $(\Omega,\mathscr{A},\mu)$. If $\sum_{n}\mu(E_n)=\infty$, then
\[
\mu\left(\limsup_{n\to\infty} E_n\right)
\geq 
\limsup_{N\to\infty} \frac{\left( \displaystyle\sum_{1\leq n\leq N} \mu(E_n)\right)^2}{\displaystyle\sum_{1\leq i\neq j\leq N} \mu(E_i\cap E_j)}.
\]
\end{lem01}





\subsection{Hausdorff dimension}
The Hausdorff dimension provides a way to distinguish between the sets of Lebesgue measure $0$. We recall the definition of the Hausdorff dimension and refer to \cite{FalconerBookFractal} for its general theory.

Given a non-empty subset $A$ of real numbers, let $|A|$ be its diameter. For $\delta > 0$, by a $\delta$-cover of $A$ we mean a countable collection of subsets of real numbers $\{B_i:i\in\Na\}$ such that 
\[
\sup_{i\in\Na} |B_i|\leq \delta
\quad\text{ and }\quad
A\subseteq \bigcup_{i\in\Na} B_i.
\]
For $s\geq 0$ and $E\subseteq \RE$ we define
\[
\clH^s_{\delta}(E)
=
\inf\left\{\sum_{i\in\Na}|B_i|^s: \{B_i:i\in\Na\} \text{ is a }\delta-\text{cover of } E\right\}.
\]
The $s$-Hausdorff measure $\clH^s$ of $E$ is the outer measure given by
\[
\clH^s(E)
:=
\lim_{\delta\to 0} \clH^s_{\delta}(E).
\]
It is well known that for each $E\subseteq \RE$ there is a number, called the Hausdorff dimension of $E$ and denoted $\hdim E$, such that for all $s\geq 0$
\[
\clH^s(E)
:=
\begin{cases}
\infty, & \text{if} \ \ s< \hdim E,\\
0, & \text{if} \ \  s> \hdim E.
\end{cases}
\]
Thus, the Hausdorff dimension of $E$ can also be defined as
\[
\hdim E
=
\inf\left\{ s\geq 0: \clH^s(E)=0\right\}.
\]
\subsection{Continued fractions}
We recall some elementary facts of continued fractions. For each real number $t$, let $[t]$ denote its integer part. Regular continued fractions can be generated with the Gauss map $T:[0,1)\to [0,1)$ given by 
\[
T(x) = 
\begin{cases}
x^{-1} - [x^{-1}], & x\neq 0, \\
0, & x\neq 0.
\end{cases}
\]
For any $x\in [0,1)\setminus \QU$ let $a_1:[0,1)\setminus \QU\to\Na$ be given by $a_1(x)= [x^{-1}]$ and for each $n\in\Na$ we define $a_n(x)=a_1(T^{n-1}(x))$.  The numbers $(a_n(x))_{n\geq 1}$ are known as the partial quotients of $x$ and they are the single sequence of natural numbers satisfying
\[
x=
[a_1(x),a_2(x),\ldots].
\]
We can extend the definition of partial quotients in an obvious manner for rational numbers. In this case, however, we obtain a finite sequence. For any non-negative integer $n$, the $n$-th convergent of $x$ is 
\[
\frac{p_n(x)}{q_n(x)}
:=
[a_1(x), \ldots, a_n(x)].
\]
An elementary result (see, for example, \cite[Proposition 1.1]{HensleyBook}) states that the sequences $(p_n(x))_{n\geq 0}$, $(q_n(x))_{n\geq 1}$ verify
\[
q_0(x)=1, \quad
q_{1}(x)=a_1(x), \quad
p_0(x)=0, \quad
p_1(x)=1
\]
and 
\[
p_{n+1}(x)=a_{n+1}(x)p_{n}(x) + p_{n-1}(x), \quad
q_{n+1}(x)=a_{n+1}(x)q_{n}(x)+q_{n-1}(x)
\;\text{ for all } n\in\Na.
\]
The terms of the sequence $(q_n(x))_{n\geq 1}$ are called continuants. Observe that the $n$-th continuant depends only on the first partial $n$ quotients of $x$. Therefore, any other number $y$ whose first $n$ partial quotients are exactly those of $x$ also has the same continuants as $x$. When we consider the partial quotients as functions on $[0,1)$, we can easily determined their inverse images. For each $n\in\Na$ and $\bfa=(a_1,\ldots, a_n)\in\Na^n$, the fundamental interval $I_n(\bfa)$ is
\[
I_n(\bfa)
:=
\left\{ x\in [0,1): a_1(x)=a_1,\ldots, a_n(x)=a_n\right\}.
\]
In the following, we denote by $q_n(\bfa)$ the $n$-th convergent of any number in $I_n(\bfa)$. We omit the argument of $q_n$ if there is no risk of ambiguity.
\begin{lem01}[{\cite[Section 12]{KhinchinCFBook}}]
There are some absolute constants such that for all $n\in\Na$ and all $\bfa\in\Na^n$ we have
\[
|I_n(\bfa)|
\asymp
\frac{1}{q_n^2(\bfa)}.
\]
\end{lem01}
For any $n,m\in\Na$, $\bfa=(a_1,\ldots, a_n)\in\Na^n$, and $\bfb=(b_1,\ldots, b_m)\in\Na^m$, we write
\[
(\bfa,\bfb)
:=
(a_1,\ldots, a_n,b_1,\ldots, b_m).
\]
\begin{lem01}[{\cite[Proposition 1.1]{HensleyBook}}]
For any $m,n\in\Na$,  $\bfa\in\Na^n$, and $\bfb\in\Na^m$, 
\[
q_{n}(\bfa)q_{m}(\bfb)
\leq
q_{n+m}(\bfa,\bfb)
\leq 
2q_{n}(\bfa)q_{m}(\bfb).
\]
\end{lem01}
\begin{lem01}[{\cite[Lemma 2.1]{Wu2006}}]
For any $n\in\Na$,  $\bfa=(a_1,\ldots, a_n)\in\Na^n$, and $k\in\{1,\ldots,n\}$, we have
\[
\frac{a_{k} + 1}{2}
\leq 
\frac{q_n(a_1,\ldots, a_n)}{q_{n-1}(a_1,\ldots,a_{k-1},a_{k+1},\ldots, a_n)}
\leq 
a_{k}+1.
\]
\end{lem01}

\begin{lem01}[\cite{FengWuLiang1997, Luc1997}]\label{Lem:Luczak}
For any constants $b,c>1$, we have 
\begin{align*}
\frac{1}{b+1}
&=
\hdim
\left\{ x\in [0,1): a_n(x)\geq b^{c^n} \text{ for all large } n\in\Na \right\}, \\
&=
\hdim
\left\{ x\in [0,1): a_n(x)\geq b^{c^n} \text{ for i.m.   } n\in\Na\right\}.
\end{align*}
\end{lem01}



\subsection{Pressure functions}
For the general theory of pressure functions, we refer the reader to the monograph by Mauldin and Urba{\'n}ski \cite{MauldinUrbanskiBook}. Pressure function have gained popularity in Diophantine approximation because their zeros are the Hausdorff dimension of certain sets determined by dynamical means. In this paper, we restrict ourselves to the dynamical system $([0,1)\setminus \QU,T)$. 

Given a function $\vphi:[0,1] \to\RE$, for $x\in [0,1]$ and $n\in\Na$, we denote by $S_n(\vphi;x)$ the ergodic sum
\[
S_n(\vphi;x)
:=
\vphi(x) + \vphi(T(x)) + \ldots + \vphi\left(T^{n-1}(x)\right).
\]
For each $n\in\Na$ the $n$--th variation of $\vphi$ is given by
\[
\Var_n(\vphi)
:=
\sup\left\{ |\vphi(x)-\vphi(y)|:I_n(x)=I_n(y)\right\}.
\]
For any non-empty subset $A\subseteq \Na$ we consider
\[
X_A
:=
\left\{ x\in [0,1): a_n(x)\in A \text{ for all } n\in\Na\right\}.
\]
The pressure function $P_{A}(T,\cdot)$ on $\vphi$ is given by
\[
P_{A}(T,\vphi)
:=
\lim_{n\to\infty} 
\frac{1}{n} 
\log\left( 
\sum_{\bfa\in A^n} \sup_{x\in I_n(\bfa)\cap X_{ A}} \exp\left( S_n(\vphi,x)\right)
\right).
\]
In this context, we call $\vphi$ a potential. When $A=\Na$ we simply write $P(T,\vphi)$. The existence of $P_{A}(T,\vphi)$ is not immediate, but it does exist under some reasonable regularity conditions. Furthermore, we can approximate $P(T,\vphi)$ by $P_{A}(T,\vphi)$ with $A$ finite.

\begin{lem01}[{\cite[Proposition 2.4]{LiWangWuXu2014}}]\label{Lem:PF:01}
Let $\vphi:[0,1]\to\RE$ be such that $\Var_1(\vphi)<\infty$ and $\Var_n(\vphi)\to 0$ as $n\to\infty$. Then, $P_{A}(T,\vphi)$ is well defined and the supremum in the definition of $P_{A}(T,\vphi)$ may be replaced by any point in $I_n(\bfa)\cap X_{A}$.
\end{lem01}
\begin{lem01}[{\cite[Proposition 2]{HanusMauldinUrbanski2002}}]\label{Lem:PF:02}
Let $\vphi:[0,1]\to\RE$ be such that $\Var_1(\vphi)<\infty$ and $\Var_n(\vphi)\to 0$ as $n\to\infty$. We have
\[
P(T,\vphi)
=
\sup \left\{ P_{A}(T,\vphi): A\subseteq \Na\;\text{\rm is finite}\right\}.
\]
\end{lem01}
Given a continuous function $f:[0,1]\to\RE$, consider the potential
\[
\vphi_s(x)
:= 
-f(s)\log B - s\log|T'(x)|.
\]
Then, $\vphi$ satisfies the hypotheses of Lemma \ref{Lem:PF:01} and of Lemma \ref{Lem:PF:02}. Since for every $x\in [0,1)$ with $n$ partial quotients we have
\[
(T^n)'(x)
=
\frac{(-1)^n}{(xq_{n-1}-p_{n-1})^2},
\]
the pressure function $P(T,\cdot)$ evaluated on $\vphi_s$ looks like this:
\[
P(T,\vphi_s)
=
\lim_{n\to\infty} 
\frac{1}{n} 
\log\left( 
\sum_{\bfa\in A^n} \frac{1}{B^{f(s)} q_{n}^{2s}(\bfa)} 
\right).
\]
In our context, for some $B>1$, we consider the potentials $\vphi_s$ given by 
\[
\vphi_s(x)
=
-g_3(s)\log B -s\log |T'(x)|.
\]
Define
\[
p_B
:=
\inf\{s\geq 0: -g_3(s)\log B -s\log |T'(x)|\}.
\]
For each $m\in\Na$ write
\begin{equation}\label{Eq:DefSmB}
s_m(B)
:=
\inf_{s\geq 0}
\sum_{(a_1,\ldots, a_m)\in\Na^m} 
\left(\frac{1}{q_m^{2s} B^{g_3(s)}}\right) \leq 1.
\end{equation}
The arguments of \cite[Section 2]{WangWu2008} allow us to conclude
\[
\lim_{m\to\infty} s_m(B)
=
s_B.
\]
Moreover, relying on Theorem \ref{TEO:TTW:HD}, \cite[Proposition 2.11]{TanTianWang2023}, and $g_3(x)\leq 3x-1$ for $x\geq \frac{2}{5}$ we may show that $p_B\geq \frac{1}{2}$.


The lower estimate for the Hausdorff dimension of sets similar to that in Theorem \ref{TEO:MAIN:HD} tends to be computationally demanding. A general strategy is to construct a Cantor set $C$ contained in our set of interest, afterwards define a probablity measure on $C$ and apply the Mass Distribution Principle \cite[Proposition 4.2]{FalconerBookFractal}. Recently, Hussain and Shulga \cite{ShulgaHussain2023Exp} computed the Hausdorff dimension of a large family of sets defined by continued fractions restrictions. For $A_0,\ldots, A_{\ell-1}>1$, define  
\[
S_{\ell}(A_0,\ldots, A_{\ell-1})
:=
\left\{ x\in [0,1): A_0^n\leq a_n(x)\leq 2A_0^n, \ldots, A_{\ell-1}^n\leq a_{n+\ell-1} (x)\leq 2A_{\ell-1}^n \text{ for i.m.   } n\in\Na\right\}.
\]
Put $\beta_{-1}=1$ and $\beta_{i}=A_i\beta_{i-1}$ for $i\in \{0,\ldots, \ell-1\}$ and for such $i$ write
\[
d_i
:=
\inf\left\{ s\geq 0: P(T,-s\log|T'| - s\log \beta_i + (1-s)\log\beta_{i-1}\leq 0\right\}.
\]
\begin{lem01}[{\cite[Theorem 1.1]{ShulgaHussain2023Exp}}]\label{Le:MumtazNikita}
$\hdim S_{\ell}(A_0,\ldots, A_{\ell-1})
=
\displaystyle\min_{0\leq i\leq \ell-1} d_i$.

\end{lem01}
\section{Series estimates}\label{Section:Series}
In this section we estimate certain series appearing in the proofs of theorems \ref{TEO:MAIN:MEASURE} and \ref{TEO:MAIN:HD}.
In all of the sums studied in this section, the variables $a_1,a_2,\ldots$ are natural numbers. For readibilty, we omit this restriction. We remind the reader that $\ell$ stands for an arbitrary natural number.

Let us start with the lemmas needed in the proof of Theorem \ref{TEO:MAIN:MEASURE}. The first one can be retrieved from the proof of \cite[Theorem 1.5]{HuangWuXu2020}.
\begin{lem01}[\cite{HuangWuXu2020}]
%For any $\ell\in\Na$ and $M>1$, we have
For any $M>1$, we have
\begin{equation}\label{Eq:Series:01}
\sum_{a_1\cdots a_{\ell}\geq M }
\frac{1}{a_1^2\cdots a_{\ell}^2} 
\asymp_{\ell}
\frac{\log^{\ell-1}(M)}{M}.
\end{equation}
\end{lem01}

We recall a result from analytic number theory (see, for example, \cite[Section 4]{HardyLittlewood1923}).

\begin{lem01}[Dirichlet-Piltz formula]
Take $k\in\Na$. For each $x\in \Na$, let $d_k(x)$ denote the number of ways we can express $x$ as a product of $k$ natural numbers. For every $M>1$, define
\[
D_k(M):= \sum_{x\leq M} d_k(x).
\]
There exist a polynomial $P_k$ of degree $k-1$ and a constant $\alpha_k<1$ such that for any $\veps>0$ we have
\[
D_k(M)
=
MP_k(\log M) + O(M^{\alpha_k+\veps})
\;\text{ as }\; M\to\infty.
\]
\end{lem01}
\begin{lem01}
For any $j\in\Na$ and large $M>1$, we have
\[
    \sum_{\substack{ab_1b_2\cdots b_j\geq M \\ b_1b_2\cdots b_jc\geq M}} 
    \frac{1}{a^2b_1^2\cdots b_j^2c^2}
    \asymp_{j}
    \frac{\log^{j-1} M}{M}, \label{Eq:01:01} 
\]
\end{lem01}
\begin{proof}
 Note that 
\begin{align*}
\sum_{\substack{ab_1\cdots b_j\geq M \\ b_1\cdots b_jc\geq M}} \frac{1}{a^2b_1^2\cdots b_j^2c^2}
&=
\sum_{b_1\cdots b_j\geq M} 
\frac{1}{b_1^2\cdots b_j^2}
\left( \sum_{a> \frac{M}{b_1\cdots b_j}} \frac{1}{a^2}\right)^2
+
\sum_{b_1\cdots b_j< M} 
\frac{1}{b_1^2\cdots b_j^2}
\left( \sum_{a> \frac{M}{b_1\cdots b_j}} \frac{1}{a^2}\right)^2
\end{align*}
When $b_1\cdots b_j\geq M$, the restriction $a \geq  \frac{M}{b_1\cdots b_j}$ is irrelevant, so
\[
\sum_{b_1\cdots b_j\geq M} 
\frac{1}{b_1^2\cdots b_j^2}
\left( \sum_{a> \frac{M}{b_1\cdots b_j}} \frac{1}{a^2}\right)^2
\asymp
\sum_{b_1\cdots b_j\geq M} 
\frac{1}{b_1^2\cdots b_j^2}
\asymp
\frac{\log^{j-1} M}{M}.
\]
For the second term, we use the Dirichlet-Piltz formula to obtain
\[
\sum_{b_1\cdots b_j< M} 
\frac{1}{b_1^2\cdots b_j^2}
\left( \sum_{a> \frac{M}{b_1\cdots b_j}} \frac{1}{a^2}\right)^2
\ll 
\sum_{b_1\cdots b_j< M} 
\frac{1}{b_1^2\cdots b_j^2}
\frac{1}{\frac{M^2}{b_1^2\cdots b_j^2}}
=
\frac{1}{M^2}
\sum_{b_1\cdots b_j< M} 
1
\asymp
\frac{\log^{j-1} M}{M}.
\]
\end{proof}

\begin{lem01}
For large $M>1$, we have
\begin{equation}\label{Eq:01:02}
\sum_{\substack{a_1a_2 b\geq M \\ b c_1c_2\geq M}} \frac{1}{a_1^2a_2^2b^2c_1^2c_2^2}
\asymp
\frac{1}{M}. 
\end{equation}
\end{lem01}
\begin{proof}
Observe that
\[
\sum_{\substack{a_1a_2 b\geq M \\ b c_1c_2\geq M}} \frac{1}{a_1^2a_2^2b^2c_1^2c_2^2}
=
\sum_{b\geq M} 
\frac{1}{b^2}
\left( \sum_{a_1 a_2 \geq \frac{M}{b}} \frac{1}{a_1^2a_2^2}\right)^2
+
\sum_{b<M} 
\frac{1}{b^2}
\left( \sum_{a_1 a_2 \geq \frac{M}{b}} \frac{1}{a_1^2a_2^2}\right)^2.
\]
The condition $a_1a_2\geq \frac{M}{b}$ is redundant when $b\geq M$, so
\[
\sum_{b\geq M} 
\frac{1}{b^2}
\left( \sum_{a_1 a_2 \geq \frac{M}{b}} \frac{1}{a_1^2a_2^2}\right)^2
\asymp
\sum_{b\geq M} 
\frac{1}{b^2}
\asymp
\frac{1}{M}.
\]
For the second term, we have
\[
\sum_{b<M} 
\frac{1}{b^2}
\left( \sum_{a_1 a_2 \geq \frac{M}{b}} \frac{1}{a_1^2a_2^2}\right)^2
\asymp
\sum_{b<M} 
\frac{1}{b^2}
\frac{\log^2 \frac{M}{b} }{\frac{M^2}{b^2}}
=
\frac{1}{M^2}
\sum_{b<M} 
\log^2 \left(\frac{M}{b}\right)
\asymp
\frac{1}{M^2}
\int_1^M
\log^2\left( \frac{x}{M}\right)
\asymp
\frac{1}{M}.
\]
\end{proof}
\begin{lem01}\label{Lem:SerEstNu}
For every large $M$ (depending on $\ell$) and any $r,j \in\Na$ such that $r+j=\ell$ we have
\begin{equation}\label{Eq:Series:02}
\sum_{\substack{a_1 \cdots a_r b_1\cdots b_j\geq M \\ b_1\cdots b_j c_1 \cdots c_r\geq M}} \frac{1}{a_1^2\cdots a_r^2 b_1^2\cdots b_j^2 c_1^2 \cdots c_r^2}
\ll_{\ell} 
\frac{\log^{2(j-1)}(M)}{M},
\end{equation}
\begin{equation}\label{Eq:Series:03}
\sum_{a_1\cdots a_{\ell}\leq M} \frac{1}{a_1\cdots a_{\ell}} 
\asymp_{\ell}
\log^{\ell}(M), 
\end{equation}
\begin{equation}  \label{Eq:Series:04}
\sum_{\substack{a_1\cdots a_{\ell}<M \\ a_2\cdots a_{\ell+1}\geq M}} \frac{1}{a_1^2\cdots a_{\ell}^2a_{\ell+1}^2 }
\asymp_{\ell}
\frac{\log^{\ell-1}M}{M}. 
\end{equation}
\end{lem01}
\begin{rema01}
The assumption $r+j=\ell$ in \eqref{Eq:Series:02} is needed for the implied constant to depend on $\ell$. Otherwise, as the proof makes it clear, this constant depends on $j$.
\end{rema01}
\begin{proof}
We start with \eqref{Eq:Series:02}. We split the sum on the left hand side of \eqref{Eq:Series:02} into two sums according to whether $b_1\cdots b_j\geq M$ or $b_1\cdots b_j< M$. In the first case, we have

\begin{align*}
\sum_{\substack{a_1 \cdots a_r b_1\cdots b_j\geq M \\ b_1\cdots b_j c_1 \cdots c_r\geq M \\b_1\cdots b_j\geq M }} \frac{1}{a_1^2\cdots a_r^2 b_1^2\cdots b_j^2 c_1^2 \cdots c_r^2}
&=
\sum_{ b_1\cdots b_j\geq M } 
\sum_{\substack{a_1 \cdots a_r b_1\cdots b_j\geq M \\ b_1\cdots b_j c_1 \cdots c_r\geq M  }} 
\frac{1}{a_1^2\cdots a_r^2 b_1^2\cdots b_j^2 c_1^2 \cdots c_r^2} \nonumber\\
&=
\sum_{ b_1\cdots b_j\geq M } 
\frac{1}{b_1^2\cdots b_j^2}
\sum_{\substack{a_1, \ldots, a_r\in\Na \\ c_1, \ldots, c_r\in\Na  }} 
\frac{1}{a_1^2\cdots a_r^2  c_1^2 \cdots c_r^2} \nonumber\\
&
\asymp \frac{\pi^{4r}}{6^{2r}}\frac{\log^{j-1}M}{M}.
\end{align*}
For the second case, we split the sum corresponding to $b_1\cdots b_j< M$ into two further sums:
\[
\sum_{\substack{a_1 \cdots a_r b_1\cdots b_j\geq M \\ b_1\cdots b_j c_1 \cdots c_r\geq M \\ \frac{M}{2}< b_1\cdots b_j< M }} \frac{1}{a_1^2\cdots a_r^2 b_1^2\cdots b_j^2 c_1^2 \cdots c_r^2}
+
\sum_{\substack{a_1 \cdots a_r b_1\cdots b_j\geq M \\ b_1\cdots b_j c_1 \cdots c_r\geq M \\ b_1\cdots b_j< \frac{M}{2} }} \frac{1}{a_1^2\cdots a_r^2 b_1^2\cdots b_j^2 c_1^2 \cdots c_r^2}
\]
Let us start with the first one. Observe that 
\[
 \frac{M}{2}< b_1\cdots b_j< M 
\quad\text{ if and only if}\quad
1< \frac{M}{b_1\cdots b_j} < 2,
\]
hence
\begin{align*}
\sum_{\substack{a_1 \cdots a_r b_1\cdots b_j\geq M \\ b_1\cdots b_j c_1 \cdots c_r\geq M \\ \frac{M}{2}< b_1\cdots b_j< M }} 
\frac{1}{a_1^2\cdots a_r^2 b_1^2\cdots b_j^2 c_1^2 \cdots c_r^2}
&=
\sum_{ \frac{M}{2}< b_1\cdots b_j< M } 
\frac{1}{b_1^2\cdots b_j^2}
\sum_{\substack{a_1 \cdots a_r\geq \frac{M}{b_1\cdots b_j} \\  c_1 \cdots c_r\geq \frac{M}{b_1\cdots b_j} }} 
\frac{1}{a_1^2\cdots a_r^2 c_1^2 \cdots c_r^2} \nonumber\\
&
=
\sum_{ \frac{M}{2}< b_1\cdots b_j< M } 
\frac{1}{b_1^2\cdots b_j^2}
\left(\left(\sum_{a\geq 1} \frac{1}{a^2}\right)^{r} -1 \right)^2 \nonumber\\
&\asymp_{\ell}
\sum_{ \frac{M}{2}< b_1\cdots b_j< M } 
\frac{1}{b_1^2\cdots b_j^2} \ll \frac{\log^{j-1}M}{M},
\end{align*}
%With more effort, we can even replace $\log^{j-1}(M)$ in the last term with $O(M)$, but it is irrelevant for our goals. 
Lastly, we have
\begin{align*}
\sum_{\substack{a_1 \cdots a_r b_1\cdots b_j\geq M \\ b_1\cdots b_j c_1 \cdots c_r\geq M \\ b_1\cdots b_j\leq \frac{M}{2}}} 
\frac{1}{a_1^2\cdots a_r^2 b_1^2\cdots b_j^2 c_1^2 \cdots c_r^2}
&=
\sum_{ b_1\cdots b_j \leq \frac{M}{2} } 
\frac{1}{b_1^2\cdots b_j^2}
\sum_{\substack{a_1 \cdots a_r\geq \frac{M}{b_1\cdots b_j} \\  c_1 \cdots c_r\geq \frac{M}{b_1\cdots b_j} }} 
\frac{1}{a_1^2\cdots a_r^2 c_1^2 \cdots c_r^2} \nonumber\\
&=
\sum_{  b_1\cdots b_j\leq  \frac{M}{2} } 
\frac{1}{b_1^2\cdots b_j^2}
\left(\sum_{ a_1 \cdots a_r\geq \frac{M}{b_1\cdots b_j} } 
\frac{1}{a_1^2\cdots a_r^2}\right)^2 \nonumber\\
&\asymp
\sum_{  b_1\cdots b_j \leq \frac{M}{2}} 
\frac{1}{b_1^2\cdots b_j^2}
\left( \frac{\log^{j-1}\left(\frac{M}{b_1\cdots b_j}\right) }{\frac{M}{b_1\cdots b_j}} \right)^2 \nonumber\\
&\asymp
\frac{1}{M^2}
\sum_{  b_1\cdots b_j \leq \frac{M}{2} } \log^{2(j-1)}\left(\frac{M}{b_1\cdots b_j}\right)
\nonumber\\
&\ll
\frac{\log^{2(j-1)}M}{M^2}
\sum_{  b_1\cdots b_j \leq \frac{M}{2} } 1
\nonumber\\
&\ll
\frac{\log^{2(j-1)}M}{M}.
\end{align*}
The last step follows from Dirichlet-Piltz formula.

The estimate \eqref{Eq:Series:03} is clear when $\ell=1$. Assume that it holds for some $\ell\in\Na$. It is easily seen that 
\begin{equation}\label{Est:Integral}
\int \frac{1}{t} \log^{\ell}\left( \frac{M}{t}\right) \md t
= 
\frac{-1}{\ell+1} \log^{\ell+1}\left( \frac{M}{t}\right) + C.
\end{equation}
We write the sum for $\ell+1$ as follows:
\[
\sum_{a_1\cdots a_{\ell}a_{\ell+1}\leq M} \frac{1}{a_1\cdots a_{\ell}} 
=
\sum_{1\leq a_{\ell+1}< \frac{M}{2}} \frac{1}{a_{\ell+1}} \sum_{a_1\cdots a_{\ell}\leq \frac{M}{a_{\ell+1}}} \frac{1}{a_1\cdots a_{\ell}} 
+
\sum_{\frac{M}{2} \leq a_{\ell+1}\leq M} \frac{1}{a_{\ell+1}} \sum_{a_1\cdots a_{\ell}\leq \frac{M}{a_{\ell+1}}} \frac{1}{a_1\cdots a_{\ell}} 
\]
We use the induction hypothesis and \eqref{Est:Integral} on the first term to get
\[
\sum_{1\leq a_{\ell+1}< \frac{M}{2}} \frac{1}{a_{\ell+1}} \sum_{a_1\cdots a_{\ell}\leq \frac{M}{a_{\ell+1}}} \frac{1}{a_1\cdots a_{\ell}} 
\asymp
\sum_{1\leq a_{\ell+1}< \frac{M}{2}} \frac{1}{a_{\ell+1}} \log^{\ell}\left( \frac{M}{a_{\ell+1}}\right)
\asymp
\log^{\ell+1}(M).
\]
The second term is bounded, because of \eqref{Est:Integral}:
\[
\sum_{\frac{M}{2} \leq a_{\ell+1}\leq M} \frac{1}{a_{\ell+1}} \sum_{a_1\cdots a_{\ell}\leq \frac{M}{a_{\ell+1}}} \frac{1}{a_1\cdots a_{\ell}} 
\asymp 
\int_{M/2}^M \frac{1}{t}\log^{\ell}\frac{M}{t} \md t
\asymp_{\ell}
1.
\]
Therefore, \eqref{Eq:Series:03} also holds for $\ell+1$.

Lastly, we prove \eqref{Eq:Series:04}. First, we rewrite the series in \eqref{Eq:Series:04} as
\[
\sum_{\substack{a_1\cdots a_{\ell}<M \\ a_2\cdots a_{\ell+1}\geq M}} \frac{1}{a_1^2\cdots  a_{\ell+1}^2 }
=
\sum_{a_2\cdots a_{\ell}\leq M}
\sum_{ a_1\leq \frac{M}{a_2\ldots a_{\ell}} } 
\sum_{a_{\ell+1} \geq \frac{M}{a_2\cdots a_{\ell}}} 
\frac{1}{a_1^2  \cdots a_{\ell+1}^2}.
\] 
We split the last series into two by conditioning on whether $1\leq a_2\ldots a_{\ell}\leq \frac{M}{2}$ or $\frac{M}{2}< a_2\cdots a_{\ell} < M$. In the first case, since $1\leq \sum_{j=1}^n j^{-2}< 3$ for all $n\in\Na$ and because of \eqref{Eq:Series:03}, we have
\begin{align*}
\sum_{a_2\cdots a_{\ell}\leq \frac{M}{2} }
\sum_{ a_1\leq \frac{M}{a_2\ldots a_{\ell}} } 
\sum_{a_{\ell+1} \geq \frac{M}{a_2\cdots a_{\ell}}} 
\frac{1}{a_1^2  \cdots a_{\ell+1}^2}
&\asymp 
\sum_{a_2\cdots a_{\ell+1}\leq \frac{M}{2} } 
\frac{1}{a_2^2\cdots a_{\ell}^2}
\sum_{ a_1\leq \frac{M}{a_2\ldots a_{\ell+1}} } 
\frac{1}{a_1^2}
\sum_{a_{\ell+2} \geq \frac{M}{a_2\cdots a_{\ell+1}}} 
\frac{1}{a_{\ell+1}^2} \\
&\asymp 
\sum_{a_2\cdots a_{\ell}\leq \frac{M}{2} }
\frac{1}{a_2^2 \cdots a_{\ell}^2}
\left( \sum_{a_{\ell+1} \geq \frac{M}{a_2\cdots a_{\ell}}} \frac{1}{a_{\ell+ 1}^2}\right)   \\
&\asymp
\frac{1}{M}
\sum_{a_2\cdots a_{\ell}\leq \frac{M}{2} }
\frac{1}{a_2 \cdots a_{\ell}} \\
&\asymp
\frac{\log^{\ell-1}M}{M}.
\end{align*}
When $\frac{M}{2}< a_2\cdots a_{\ell} < M$, the condition $a_1\leq \frac{M}{a_2\ldots a_{\ell}}$ reduces to $a_1=1$ and $a_{\ell+1} \geq \frac{M}{a_2\cdots a_{\ell}}$ reduces to $a_{\ell+1}\geq 2$, then
\begin{align*}
\sum_{\frac{M}{2}\leq a_2\cdots a_{\ell}\leq M }
\sum_{ a_1\leq \frac{M}{a_2\ldots a_{\ell}} } 
\sum_{a_{\ell+2} \geq \frac{M}{a_1\cdots a_{\ell}}} 
\frac{1}{a_1^2 \cdots a_{\ell+1}^2}
&\asymp
\sum_{\frac{M}{2}\leq a_2\cdots a_{\ell}\leq M }
\frac{1}{a_2^2 \cdots a_{\ell}^2} \\
&\ll
\sum_{a_2\cdots a_{\ell}\leq \frac{M}{2} }
\frac{1}{a_2^2 \cdots a_{\ell}^2} 
\asymp
\frac{\log^{\ell-2}M}{M}.
\end{align*}
This proves \eqref{Eq:Series:04}.
\end{proof}
Theorem \ref{TEO:TTW:HD} requires similar estimates where the exponents may take non-integer values.

 \begin{lem01}[{\cite[Lemma 4.2]{HuangWuXu2020}}]
For $\beta>1$, $0<s<1$, and $n\in\Na$, we have
\begin{equation}\label{Eq:Series:05}
\sum_{ a_1\cdots a_{\ell} \leq \beta^n}
\left(\frac{1}{a_1\cdots a_{\ell}} \right)^s
\asymp
\frac{(\log \beta^n)^{{\ell}-1}}{({\ell}-1)!} \beta^{n(1-s)}.  
\end{equation}
\end{lem01}
\begin{lem01}
For any $M>1$ and any $t>1$ we have 
\begin{equation}\label{Eq:Series:06}
\sum_{ab\geq M} \frac{1}{a^tb^t}
\asymp 
\frac{M^{1-t}}{t-1} \log M,
\end{equation}
\begin{equation}\label{Eq:Series:07}
\sum_{abc\geq M} \frac{1}{a^tb^tc^t}
\ll
\left(\frac{1}{t-1} + \log M\right)M^{1-t}.
\end{equation}
\end{lem01}
\begin{proof}
The estimate \eqref{Eq:Series:06} is shown as \eqref{Eq:Series:04} and \eqref{Eq:Series:07} follows from \eqref{Eq:Series:06}.
\end{proof}


\section{Proof of Theorem \ref{TEO:MAIN:MEASURE}}\label{Section:ProofMeasure}
We show a slightly different statement for the convergence part which yields the convergence part in Theorem \ref{TEO:TZ:MEAS}. To this end, for $\ell \in\Na_{\geq 3}$, $n\in\Na_{\geq \ell}$, and $k\in \{1,\ldots, n-1\}$, define
\[
A_{n,k}(\vphi)
:=
\left\{x\in [0,1) : 
\begin{matrix}
a_k(x)\cdots a_{k+\ell - 1}(x) \geq \vphi(n), \\
a_n(x)\cdots a_{n+\ell - 1}(x) \geq \vphi(n)
\end{matrix}
 \right\}
\]
and write
\[
A_n(\vphi)
:=
\bigcup_{k=1}^{n-1} A_{n,k}(\vphi).
\]
As in \cite{TanZhou2024}, we may check that the monotonocity of $\vphi$ implies
\[
\clF_{\ell}(\vphi)
=
\limsup_{n\to\infty} A_n(\vphi).
\]
\subsection{Convergence case}
Take $n\in\Na_{\geq \ell}$. We show that 
\begin{equation}\label{Eq:ConverGen}
\frac{n\log^{2(\ell-1)}\vphi(n)}{\vphi^2(n)} + \frac{\log^{2(\ell-2)}\vphi(n)}{\vphi(n)}<\infty
\quad\text{ implies }\quad
\leb(\clF_{\ell}(\vphi))=0.
\end{equation}
By the end of the proof, we say how to modify the argument to obtain the convergence case in Theorem \ref{TEO:MAIN:MEASURE}.
We estimate $\leb(A_{n,k})$ by considering two cases according to whether the restricted blocks overlap or not.

\begin{lem01}
Let $k\in\{1,\ldots, n-1\}$.
\begin{enumerate}[1.]
\item If $1\leq k \leq n-\ell$, then 
\[
\leb( A_{n,k})
\asymp_{\ell} 
\frac{\log^{2(\ell-1)}\vphi(n)}{\vphi^2(n)}.
\]
\item If $n-\ell +1 \leq k\leq n- 1$, then 
\[
\leb( A_{n,k})
\ll_{\ell} 
\frac{\log^{2(\ell+k-n-1)}\vphi(n)}{\vphi(n)}.
\]
\end{enumerate}
\end{lem01}
\begin{proof}
\begin{enumerate}[1.]
\item Assume that $1\leq k\leq n-\ell$, so $k+\ell-1 \leq n-1$. Let $\bfb=(b_1,\ldots, b_{n+\ell-1})\in\Na^{n+\ell-1}$ such that $A_{n,k}\cap I(\bfb)\neq \vac$. Clearly $I(\bfb)\subseteq A_{n,k}$. Write
\[
\bfb'=(b_1,\ldots, b_{k-1},b_{k+\ell},\ldots, b_{n-1}),
\]
that is, $\bfb'$ is $\bfb$ without the blocks $(b_{k},\ldots, b_{k+\ell-1})$ and $(b_n,\ldots, b_{n+\ell-1})$, then
\[
\leb(I_{n+\ell-1}(\bfb))
\asymp_{\ell}
\leb(I_{n-\ell-1}(\bfb')) \, 
\frac{1}{(b_{k}^2\cdots b_{k+\ell-1}^2) (b_{n}^2\cdots b_{n+\ell-1}^2)}.
\]
Therefore, we have
\begin{align*}
\leb(A_{n,k})
&=
\sum_{\bfb\in \Na^{n+\ell-1}} \leb\left(A_{n,k}\cap I_{n+\ell-1}(\bfb)\right)
\nonumber\\
&=
\sum_{\substack{\bfb\in \Na^{n+\ell-1}\\ b_{k}\cdots b_{k+\ell-1}\geq \vphi(n)\\ b_{n}\cdots b_{n+\ell-1}\geq \vphi(n) }} 
\leb\left(I_{n+\ell-1}(\bfb)\right) \nonumber\\
&\asymp
\sum_{\substack{\bfb'\in \Na^{n-\ell-1}\\ b_{k}\cdots b_{k+\ell-1}\geq \vphi(n)\\ b_{n}\cdots b_{n+\ell-1}\geq \vphi(n) }} 
\leb(I_{n-\ell-1}(\bfb')) \, 
\frac{1}{(b_{k}^2\cdots b_{k+\ell-1}^2) (b_{n}^2\cdots b_{n+\ell-1}^2)}.
\nonumber\\
&=
\left(\sum_{\bfb'\in \Na^{n-\ell-1}} 
\leb(I_{n-\ell-1}(\bfb') ) \right)
\left(\sum_{b_{k}\cdots b_{k+\ell-1}\geq \vphi(n) } 
\frac{1}{b_{k}^2\cdots b_{k+\ell-1}^2}\right)
\left(\sum_{b_{n}\cdots b_{n+\ell-1}\geq \vphi(n) } 
\frac{1}{b_{n}^2\cdots b_{n+\ell-1}^2}\right) \nonumber\\
&=
\left(\sum_{b_{1}\cdots b_{\ell}\geq \vphi(n) } 
\frac{1}{b_{1}^2\cdots b_{\ell}^2}\right)^2 
\asymp
\frac{\log^{2(\ell-1)}\vphi(n)}{\vphi^2(n)}. 
\end{align*}
The last estimate follows from \eqref{Eq:Series:01}.
\item Assume that $n-\ell +1 \leq k\leq n- 1$ and write $j=\ell+k-n$. Then, $j$ is the length of the overlap of the blocks appearing in the definition of $A_{n,k}$. We may argue as in the previous case and apply \eqref{Eq:Series:02} rather than \eqref{Eq:Series:01} to conclude
\[
\leb(A_{n,k})
\asymp_{\ell}
\sum_{\substack{b_{k}\cdots b_{k+\ell-1}\geq \vphi(n)\\b_{n}\cdots b_{n+\ell-1}\geq \vphi(n)}}
\frac{1}{b_{k}^2\cdots b_{k+\ell-1}^2b_{k+\ell}^2 \cdots b_{n+\ell-1}^2}
\ll 
\frac{\log^{2(j-1)} \vphi(n)}{\vphi(n)}.
\]
\end{enumerate}
\end{proof}
%\begin{proof}[Proof of Theorem \ref{TEO:MAIN:MEASURE}. Convergence case.]
\begin{proof}[Proof of \eqref{Eq:ConverGen} ]
For any $n\in\Na_{\geq \ell}$ we have
\begin{align*}
\leb(A_n)
&\leq 
\sum_{k=1}^{n} \leb(A_{n,k}) \nonumber\\
&=
\sum_{k=1}^{n-\ell} \leb(A_{n,k}) + \sum_{k=n-\ell+1}^{n-1} \leb(A_{n,k}) \nonumber\\
&\ll
(n-\ell)\frac{\log^{2(\ell-1)}\vphi(n)}{\vphi^2(n)} + \frac{1}{\vphi(n)}\sum_{k=n-\ell+1}^{n-1} \log^{2(k+\ell-n-1)}\vphi(n) \nonumber\\
&\asymp
\frac{n\log^{2(\ell-1)}\vphi(n)}{\vphi^2(n)} + \frac{\log^{2(\ell-2)}\vphi(n)}{\vphi(n)}.
\end{align*}
The result follows from the convergence part of the Borel-Cantelli Lemma.
\end{proof}
For the convergence part in Theorem \ref{TEO:MAIN:MEASURE}, set $\ell=3$ and use \eqref{Eq:01:02} instead of \eqref{Eq:Series:02} in the previous argument.
\subsection{Divergence case}
We may impose two additional assumptions without losing any generality. First, we suppose that 
\[
\lim_{n\to\infty} \vphi(n)=\infty.
\]
%If this was not the case, then every irrational number with unbounded partial quotients would belong to $\clF_{\ell}(\vphi)$ and, thus, $\leb(\clF_{\ell}(\vphi))=1$ by the Borel-Bernstein Theorem. 
If this was not the case, then every irrational number with unbounded partial quotients would belong to $\clF_{3}(\vphi)$ and, thus, $\leb(\clF_{3}(\vphi))=1$ by the Borel-Bernstein Theorem. 
Second, we may assume that \footnote{There seems to be a gap in the proof of this assertion for $\ell=2$ in \cite{TanZhou2024}. Consider $0<\alpha<1$ and $\vphi(n)=n(\log n)^{\alpha}$. Then, $\sum_n \vphi(n)^{-1}=\infty$ and $\vphi(n)< n\log \vphi(n)$ for all large $n$. According to \cite[Section 3.2]{TanZhou2024}, we should consider the function $\psi$ given by $\psi(n)=\max\{\vphi(n), n\log\vphi(n)\}$. However, for all large $n$ we have $n\log \psi(n) = n\log (n\log \vphi(n)) > n\log \vphi(n) = \psi(n)$. This minor gap does not affect validity of that paper at all.} 
%\begin{equation}\label{Eq:WLOG}
%\varphi(n) \geq n \log^{\ell-1}\varphi(n)
%\quad
%\text{ for all }n\in \Na.
%\end{equation}
\begin{equation}\label{Eq:WLOG}
\varphi(n) \geq n \log^{2}\varphi(n)
\quad
\text{ for all }n\in \Na.
\end{equation}
Certainly, for each $n\in\Na$ let $x_n$ be the solution $x$ of the equation
%\[     \frac{x}{\log^{\ell-1}x} = n.  \]
\[
    \frac{x}{\log^{2}x} = n. 
\]
Then, $(x_n)_{n\geq 1}$ tends to infinity and it is eventually strictly increasing. Let $\psi: \mathbb{N} \to \mathbb{N}$ be given by 
\[
    \psi(n) = \max\{\varphi(n),x_n\}
    \quad
    \text{ for all }n\in\Na.
\]
We can trivially modify at most finitely many terms of $\psi$ to ensure that it is non-decreasing. Furthermore, since $\psi \geq \varphi$, we have $\clE_{3}(\psi) \subseteq \clE_{3}(\varphi)$ and $\psi(n) \to \infty$ as $n\to\infty$. 
It is also clear that $\psi(n) \geq n\log^{2} \psi(n)$ for all $n\in\Na$, because $x \mapsto \frac{x}{\log^{2}x}$ is strictly increasing for large $x$. Lastly, let us show that
%We can trivially modify at most finitely many terms of $\psi$ to ensure that it is non-decreasing. Furthermore, since $\psi \geq \varphi$, we have $\clE_{\ell}(\psi) \subseteq \clE_{\ell}(\varphi)$ and $\psi(n) \to \infty$ as $n\to\infty$.
%It is also clear that $\psi(n) \geq n\log^{\ell-1} \psi(n)$ for all $n\in\Na$, because $x \mapsto \frac{x}{\log^{\ell-1}x}$ is strictly increasing for large $x$. Lastly, let us show that
%\[     \sum_{n=1}^\infty n \left(\frac{\log^{\ell-1}\psi(n)}{\psi(n)}\right)^2 + \frac{\log^{\ell-2}\psi(n)}{\psi(n)} = \infty. \]
\[
    \sum_{n=1}^\infty n \left(\frac{\log^{2}\psi(n)}{\psi(n)}\right)^2 + \frac{\log \psi(n)}{\psi(n)} = \infty.
\]
The divergence is obvious if $\varphi(n)>x_n$ for all but finitely many $n \in \mathbb{N}$. If this is not the case, then there is a strictly increasing sequence of natural numbers $(m_k)_{k\geq 1}$ such that
\[
m_{k+1}>2m_k \text{ and } \varphi(m_k) < x_{m_k}
\;\text{ for all }\;k\in\Na.
\]
Hence,
\begin{comment}
\begin{align*}
    \sum_{n=1}^\infty n \left(\frac{\log^{\ell-1}\psi(n)}{\psi(n)}\right)^2 + \frac{\log^{\ell-2}\psi(n)}{\psi(n)} 
    &\geq 
    \sum_{k=1}^\infty 
    \sum_{\left[\frac{m_{k+1}}{2}\right] \leq n \leq m_{k+1}} 
     n \left(\frac{\log^{\ell-1}\psi(n)}{\psi(n)}\right)^2\\
    &\geq 
    \sum_{k=1}^\infty \frac{m_{k+1}}{2} 
    \sum_{\left[\frac{m_{k+1}}{2}\right] \leq n \leq m_{k+1}} 
    \left(\frac{\log^{\ell-1}\psi(n)}{\psi(n)}\right)^2\\
    &
    \geq 
    \sum_{k=1}^\infty \frac{m_{k+1}}{2} 
    \sum_{\left[\frac{m_{k+1}}{2}\right] \leq n \leq m_{k+1}} 
    \left(\frac{\log^{\ell-1}x_{m_{k+1}}}{x_{m_{k+1}}}\right)^2\\
    &= 
    \sum_{k=1}^\infty \frac{m_{k+1}}{2} 
    \sum_{\left[\frac{m_{k+1}}{2}\right] \leq n \leq m_{k+1}} 
    \left(\frac{1}{m_{k+1}}\right)^2 \\
    &\geq 
    \sum_{k=1}^\infty \frac{1}{4}=\infty.
\end{align*}
\end{comment}
\begin{align*}
    \sum_{n=1}^\infty n \left(\frac{\log^{2}\psi(n)}{\psi(n)}\right)^2 + \frac{\log \psi(n)}{\psi(n)} 
    &\geq 
    \sum_{k=1}^\infty 
    \sum_{\left[\frac{m_{k+1}}{2}\right] \leq n \leq m_{k+1}} 
     n \left(\frac{\log^{2}\psi(n)}{\psi(n)}\right)^2\\
    &\geq 
    \sum_{k=1}^\infty \frac{m_{k+1}}{2} 
    \sum_{\left[\frac{m_{k+1}}{2}\right] \leq n \leq m_{k+1}} 
    \left(\frac{\log^{2}\psi(n)}{\psi(n)}\right)^2\\
    &
    \geq 
    \sum_{k=1}^\infty \frac{m_{k+1}}{2} 
    \sum_{\left[\frac{m_{k+1}}{2}\right] \leq n \leq m_{k+1}} 
    \left(\frac{\log^{2}x_{m_{k+1}}}{x_{m_{k+1}}}\right)^2\\
    &= 
    \sum_{k=1}^\infty \frac{m_{k+1}}{2} 
    \sum_{\left[\frac{m_{k+1}}{2}\right] \leq n \leq m_{k+1}} 
    \left(\frac{1}{m_{k+1}}\right)^2 \\
    &\geq 
    \sum_{k=1}^\infty \frac{1}{4}=\infty.
\end{align*}
Take $t\in\Na$, $\bfb=(b_1,\ldots, b_t)\in \Na^t$, write $I_t=I_t(\bfb)$, and let $M=M(\bfb)\in\Na_{\geq t+1}$ be such that $\vphi(M)>\left(\max\{b_1,\ldots, b_t\}\right)^{3}$.
\begin{lem01}\label{Lem:Divergence:01}
If $n\in\Na_{\geq M}$, then
\[
\leb\left( A_n\cap I_t\right)
\gg 
\leb(I_t)
\left( 
\frac{n\log^{4}\vphi(n)}{\vphi^2(n)} 
+  
\frac{\log\vphi(n)}{\vphi(n)} 
\right).
\]
\end{lem01}
\begin{proof}
{
Take $n\in\Na_{\geq M}$. Define the set
\[
B_{n}
:=
\left\{x\in I_t: 
\begin{matrix}
a_{n-1}(x)a_{n}(x)a_{n+1}(x) \geq \vphi(n), \\ 
a_{n}(x)a_{n+1}(x)a_{n+2}(x) \geq \vphi(n)
\end{matrix}
\right\},
\]
%then
%\[\leb (B_n) \asymp \frac{\log\vphi(n)}{\vphi(n)}\leb(I_t).\]
For each $i\in\Na$ such that $M\leq i<n-1$ and $i\equiv n-1 \pmod 3$, let $B_i$ be the set of those $x\in I_t$ satisfying the next properties:
\begin{enumerate}[1.]
\item $a_{i}(x)a_{i+1}(x)a_{i+2}(x) \geq \vphi(n)$, 
\item For each $k\in\Na$ such that $0< k < \frac{n-1-i}{3}$ we have $a_{i+3k}(x)a_{i+1+3k}(x)a_{i+2+3k}(x) < \vphi(n)$.
 \item $a_{n-1}(x)a_{n}(x)a_{n+1}(x) <\vphi(n)$ and $a_{n}(x)a_{n+1}(x)a_{n+2}(x) \geq \vphi(n)$.
\end{enumerate}
Note that we are imposing restrictions on non-overlapping contiguous blocks of length $3$. Clearly, we have 
\begin{equation}\label{Eq:Conten}
B_n\cup \bigcup_{\substack{M\leq  i<n-1\\ i\equiv n-1 (3)}} B_i
\subseteq
A_n\cap I_t,
\end{equation}
where the union is disjoint. From \eqref{Eq:Series:02} we obtain
\begin{equation}\label{Eq:BndBn}
\leb (B_n)
\asymp
\frac{\log\vphi(n)}{\vphi(n)}
\leb(I_t).
\end{equation}
Take $i$ such that $i\equiv n-1\pmod 3$ and $M\leq  i<n-1$. As in the convergence case, in order to estimate $\leb(B_i)$ it suffices to estimate the series
\[
\sum_{a_ia_{i+1} a_{i+2}\geq \vphi(n)} \,
\sum_{\substack{a_{i+3k}a_{i+1+3k}a_{i+2+3k } < \vphi(n) \\ 0<k<\frac{n-1-i}{3}}}\,
\sum_{\substack{a_{n-1}a_{n}a_{n+1} < \vphi(n) \\ a_{n}a_{n+1}a_{n+2} \geq  \vphi(n)}} 
\frac{1}{a_{i}^2a_{i+1}^2a_{i+2}^2 \cdots a_{n-1}^2a_{n}^2a_{n+1}^2}.
\]
}
From \eqref{Eq:Series:01} and \eqref{Eq:Series:04} we obtain
\[
\sum_{a_{i}a_{i+1}a_{i+2} \geq \vphi(n)} 
\frac{1}{a_{i}^2a_{i+1}^2a_{i+2}^2 }
\;
\sum_{\substack{a_{n-1}a_n  a_{n+1}< \vphi(n) \\ a_{n}a_{n+1}a_{n+2} \geq  \vphi(n)}} 
\frac{1}{a_{n-1}^2a_{n}^2a_{n+1}^2}
\asymp 
\frac{\log^{4} \vphi(n) }{\vphi^2(n)}.
\]
Also, for every integer $k$ with $0<k<\frac{n-1 - i}{3}$ we have
\[
\sum_{a_{i+3k}a_{i+1+3k}a_{i+2 + 3k}  < \vphi(n)}
\frac{1}{a_{i+3k}^2a_{i+1+3k}^2a_{i+2+3k}^2 }
\gg 1- \frac{\log^{2}\vphi(n)}{2\vphi(n)}
\geq 1- \frac{1}{2n}.
\]
Hence, by \eqref{Eq:Conten}, we arrive at 
\[
\leb(B_i)
\gg 
\leb(I_t)
\frac{\log^{4}\vphi(n)}{\vphi(n)}\left(1-\frac{1}{2n}\right)^{n-i}
\]
and we conclude
\begin{equation}\label{Eq:BndBn02}
\leb\left(\bigcup_{\substack{M\leq i<n-1\\ i\equiv n-1 (3)}} B_i\right)
\gg
\leb(I_t)
\frac{n\log^{4}\vphi(n)}{\vphi^2(n)}.
\end{equation}
The lemma follows from \eqref{Eq:Conten}, \eqref{Eq:BndBn}, and \eqref{Eq:BndBn02}.
\end{proof} 

\begin{lem01}\label{Lem:Divergence:02}
For large $N\in\Na$ we have
\begin{equation}\label{Eq:Lem:4.3}
\sum_{M\leq m< n \leq N} \leb\left(A_m\cap A_n\cap I_t\right)
\ll
\sum_{n=M}^N \leb\left( A_n\cap I_t \right)
+
\leb(I_t)^{-1} \left(\sum_{n=M}^N \leb\left( A_n\cap I_t \right) \right)^2.
\end{equation}
\end{lem01}
\begin{proof}
Consider the index set
\[
\clI:=\{(m,n)\in\Na^2: M\leq m< n \leq N\}.
\]
We partition $\clI$ into
%\[\clI_1:=\{(m,n)\in \clI: M\leq m \leq n - (2\ell-1)\}, \quad\clI_2:=\{(m,n)\in \clI: n-2\ell+2\leq m \leq n-1 \}  \}\]
\[
\clI_1:=\{(m,n)\in \clI: M\leq m \leq n - 5\}, 
\quad
\clI_2:=\{(m,n)\in \clI: n - 4\leq m \leq n-1 \}  \}
\]
and write
\[
S_1:=\sum_{(m,n)\in \clI_1} \leb\left(A_m\cap A_n\cap I_t\right)
\;\text{ and }\;
S_2:=\sum_{(m,n)\in \clI_2} \leb\left(A_m\cap A_n\cap I_t\right).
\]
It is clear that 
\[
S_2
\leq 
6
\sum_{n=M}^N \leb\left(A_n\cap I_t\right).
\]
Now we give an upper bound for $S_1$. For each $(m,n)\in \clI_1$ define the sets $U_1(m,n)$, $U_2(m,n)$, $U_3(m,n)$ as follows:
\begin{enumerate}[1.]
\item Let $U_1(m,n)$ be the set of irrational numbers $x\in I_t$ such that
\begin{itemize}
\item[•] $a_{m}(x)a_{m+1}(x) a_{m + 2}(x)\geq \vphi(m)$,
\item[•] For some $k<m$ we have $a_{k}(x)a_{k+1}(x) a_{k+2}(x)\geq \vphi(n)$,
\item[•] $a_{n}(x)a_{n+1}(x)a_{n+2}(x) \geq \vphi(n)$.
\end{itemize}

\item Let $U_2(m,n)$ be the set of irrational numbers $x\in I_t$ such that
\begin{itemize}
\item[•] For some $j<m$ we have $a_{j}(x)a_{j+1}(x)a_{j+2}(x) \geq \vphi(m)$,
\item[•] $a_{m}(x) a_{m + 1}(x) a_{m+2}(x) \geq \vphi(n)$,
\item[•] $a_n(x) a_{n + 1}(x)a_{n + 2}(x) \geq \vphi(n)$.
\end{itemize}

\item Let $U_3(m,n)$ be the set of irrational numbers $x \in I_t$ such that
\begin{itemize}
\item[•] For some $j<m$ we have $a_{j}(x)a_{j+1}(x)a_{j+2}(x)\geq \vphi(m)$,
\item[•] $a_{m}(x)a_{m+1}(x) a_{m + 2}(x)\geq \vphi(m)$,
\item[•] For some $m<k<n$ we have $a_{k}(x)a_{k+1}(x)a_{k+2}(x) \geq \vphi(n)$,
\item[•] $a_n(x) a_{n + 1}(x)a_{n + 2}(x)\geq \vphi(n)$.
\end{itemize}
\end{enumerate}
We claim that
\begin{equation}\label{Eq:EstU1}
\leb(U_1(m,n))
\ll
\leb(I_t) \frac{\log^{4}\vphi(n)}{\vphi^2(n)},    
\end{equation}
\begin{equation}\label{Eq:EstU2}
\leb(U_2(m,n))
\ll
\leb(I_t) \frac{\log^{4}\vphi(n)}{\vphi^2(n)}, 
\end{equation}
\begin{align}\label{Eq:EstU3}
\leb(U_3(m,n))
&\ll
\leb(I_t)
\frac{\log^{4}\vphi(n)}{\vphi^2(n)}  +\nonumber\\
&\quad\; + 
\leb(I_t)\left( \frac{n\log^{4}  \vphi(n)}{\vphi^2(n)} 
+ \frac{\log \vphi(n)}{\vphi(n)} \right)
\left( \frac{m\log^{4}  \vphi(m)}{\vphi^2(m)} + \frac{\log \vphi(m)}{\vphi(m)} \right).
\end{align}

Let us show \eqref{Eq:EstU1}. When $M\leq k \leq m-3$, we apply \eqref{Eq:Series:01} to get
\[
\sum_{\substack{a_ka_{k+1}a_{k+2} \geq \vphi(n) \\ a_ma_{m+1}a_{m+2} \geq \vphi(m) \\ a_na_{n+1}a_{n+2} \geq \vphi(n)} }
\frac{1}{a_k^2a_{k+1}^2 a_{k+2}^2 a_m^2a_{m+1}^2a_{m+2}^2 a_{n}^2a_{n+1}^2a_{n+2}^2 }
\asymp
\frac{\log\vphi(m)}{\vphi(m)}
\frac{\log^{4}\vphi(n)}{\vphi^2(n)}.
\]
For $m-2\leq k\leq m-1$ we have
\begin{align*}
\sum_{\substack{ a_ka_{k+1}a_{k+2} \geq \vphi(n)\\ a_ma_{m+1}a_{m+2} \geq \vphi(m)\\ a_{n} a_{n+1} a_{n+2}  \geq \vphi(n)}}
\frac{1}{a_{k}^2 \cdots a_{m + 2}^2 a_{n}^2a_{n+1}^2a_{n+2}^2} 
&\leq
\sum_{\substack{ a_ka_{k+1}a_{k+2}\geq \vphi(n)\\ a_{n} a_{n+1} a_{n+2}\geq \vphi(n)}}
\frac{1}{a_{k}^2 \cdots a_{m + 2}^2 a_{n}^2a_{n+1}^2a_{n+2}^2} \\
&\ll
\frac{\log^{4}\vphi(n)}{\vphi^2(n)}.
\end{align*}
Then, the above inequalities along with \eqref{Eq:WLOG} yield
\begin{align*}
\leb(U_1(m,n))
&\leq 
m\frac{\log^{2}\vphi(m)}{\vphi(m)}
\frac{\log^{4}\vphi(n)}{\vphi^2(n)}
+
3 \frac{\log^{4}\vphi(n)}{\vphi^2(n)} \\
&\ll
\frac{\log^{4}\vphi(n)}{\vphi^2(n)}.
\end{align*}
This shows \eqref{Eq:EstU1}. Note that we have used that for the values of $k$ considered the blocks $(a_k,a_{k+1}, a_{k+2})$ and $(a_n,a_{n+1}, a_{n+2})$ do not overlap. Here and in the following, whenever we assert something about blocks of a sequence of natural numbers without quantifiers, it should be understood that the assertion holds for every sequence of natural numbers.

We consider two different series to estimate $\leb(U_2(m,n))$. First, if $M\leq j\leq m- 3$ then the blocks $(a_j,a_{j+1}, a_{j+2})$ and $(a_m,a_{m+1},a_{m+2})$ do not overlap and
\begin{align*}
\sum_{j=M}^{m-\ell} 
\sum_{\substack{ a_j a_{j+1} a_{j+2}\geq \vphi(m) \\ a_ma_{m+1} a_{m+2}\geq \vphi(n)}}
\frac{1}{a_j^2 a_{j+1}^2 a_{j+2}^2 a_m^2a_{m+1}^2 a_{m+2}^2}
&\ll
m \frac{\log^{2}\vphi(m)}{\vphi(m)}\frac{\log^{2}\vphi(n)}{\vphi(n)} \\
&\leq
\frac{\log^{2}\vphi(n)}{\vphi(n)}.
\end{align*}
The product $a_j^2 \cdots a_{m+2}^2$ in the series to come should be interpreted as follows: if $j+2\leq m$, then $a_j^2 \cdots a_{m+2}^2$ is the product of the products of the blocks $(a_j,a_{j+1}, a_{j+2})$ and $(a_m, a_{m+1}, a_{m+2})$; if $j+2> m$, then $a_j^2 \cdots a_{m+2}^2$ is the product of the block $(a_{j},a_{j+1}, \ldots, a_{m+2})$. We use a similar convention for other combinations of the parameters $j,k,m,n$.

When $m - 2 \leq j \leq m-1$, the blocks $(a_j,a_{j+1}, a_{j+2})$ and $(a_m,a_{m+1}, a_{m+2})$ overlap and
\begin{align*}
\sum_{j=m - \ell +1}^{m-1} 
\sum_{\substack{ a_j a_{j+1} a_{j+2}\geq \vphi(m) \\ a_m a_{m+1} a_{m+2}\geq \vphi(n)}}
\frac{1}{a_j^2 \cdots a_{m+2}^2}
&\leq 
\sum_{j=m - 2}^{m-1} 
\sum_{a_m a_{m+1} a_{m+2}\geq \vphi(n)}
\frac{1}{a_j^2 \cdots a_{m+2}^2}\\
&\ll
\frac{\log^{2}\vphi(n)}{\vphi(n)}.
\end{align*}
We may now argue as in the convergence case recalling the restriction $a_n a_{n+1} a_{n+2}\geq \vphi(n)$ to conclude \eqref{Eq:EstU2}.

Next, we show \eqref{Eq:EstU3}. First, we consider the values of $k$ for which the block $(a_k,a_{k+1}, a_{k+2})$ does not overlap neither $(a_n,a_{n+1}, a_{n+2})$ nor $(a_m,a_{m+1}, a_{m+2})$. In other words, we assume $m+3 \leq k \leq n-3$. This leads us to the series
\[
\sum_{j=M}^{m-1}
\sum_{k=m+3}^{n-3}
\sum_{\substack{a_j a_{j+1} a_{j+2} \geq \vphi(m)\\ a_m a_{m+1} a_{m+2}\geq \vphi(m)\\ a_k a_{k+1} a_{k+2}\geq \vphi(n)\\a_na_{n+1} a_{n+2}\geq \vphi(n)}}
\frac{1}{a_j^2 \cdots a_{m+\ell-1}^2 a_k^2 a_{k+1}^2 a_{k+2}^2 a_n^2a_{n+1}^2 a_{n+2}^2},
\]
which is asymptotically equivalent to
\begin{align*}
\left( \frac{\log^{2}\vphi(n)}{\vphi(n)}\right)^2
&\sum_{j=M}^{m-1}
\sum_{\substack{a_j a_{j+1} a_{j+2}\geq \vphi(m)\\ a_m a_{m+1} a_{m+2}\geq \vphi(m)}}
\frac{1}{a_j^2 \cdots a_{m+2}^2}
\ll \\
&\ll
\left( \frac{\log^{2}\vphi(n)}{\vphi(n)}\right)^2
\left( m \frac{\log^{4}\vphi(m)}{\vphi^2(m)} + \frac{\log\vphi(m)}{\vphi(m)}\right).
\end{align*}
%{\color{red}
The last estimate is obtained as in the convergence case. 

Now suppose that $n-2\leq k \leq n-1$. This means that $(a_k,a_{k+1}, a_{k+2})$ and $(a_n,a_{n+1},,a_{n+2} )$ overlap but $(a_k,a_{k+1}, a_{k + 2})$ and $(a_m,a_{m+1}, a_{m + 2})$ do not (by the definition of $\clI_1$), hence
\begin{align*}
\sum_{j=M}^{m-1}
&\sum_{k=n - 2}^{n-1}
\sum_{\substack{a_j a_{j+1} a_{j+2}\geq \vphi(m)\\ a_m a_{m+1} a_{m+2}\geq \vphi(m)\\ a_k a_{k+1} a_{k+2}\geq \vphi(n)\\ a_n a_{n+1} a_{n+2}\geq \vphi(n)}}
\frac{1}{a_j^2 \cdots a_{m + 2}^2a_k^2\cdots a_{n + 2}^2} \ll \\
&\ll
\frac{\log \vphi(n)}{\vphi(n)}
\sum_{j=M}^{m-1}
\sum_{\substack{a_ja_{j + 1} a_{j + 2}\geq \vphi(m) \\a_m a_{m+1} a_{m+2}\geq \vphi(m) }}
\frac{1}{a_j^2 \cdots a_{m+2}^2} \\
&\ll
\frac{\log \vphi(n)}{\vphi(n)}
\left( m \frac{\log^{4}\vphi(m)}{\vphi(m)} + \frac{\log \vphi(m)}{\vphi(m)}\right).
\end{align*}
%}

Lastly, we deal with the values of $k$ for which the blocks $(a_m,a_{m+1}, a_{m+2})$ and $(a_k, a_{k+1}, a_{k+2})$ overlap. This happens if and only if $k\leq m+2$. In this case, $k+ 2 \leq m+4<n$, hence $(a_k, a_{k+1}, a_{k+2})$ and $(a_n, a_{n+1}, a_{n+2})$ do not overlap. For any $j\in \{M, \ldots, m-1\}$, $k\in \{m+1,m+2\}$ define
\[
D_{j,k}
:=
\left\{
x\in I_t:
\begin{matrix}
a_j(x) a_{j+1}(x) a_{j+2}(x)\geq \vphi(m), & a_m(x) a_{m+1}(x) a_{m+2}(x)\geq \vphi(m), \\
a_k(x) a_{k+1}(x) a_{k+2}(x)\geq \vphi(n), & a_n(x) a_{n+1}(x) a_{n+2}(x)\geq \vphi(n)
\end{matrix}
\right\}.
\]
For such $k$ we have
\[
\bigcup_{j=M}^{m-1} D_{j,k}
\subseteq
\left\{
x\in I_t:
\begin{matrix}
a_k(x) a_{k+1}(x) a_{k+2}(x)\geq \vphi(n),\\  
a_n(x) a_{n+1}(x) a_{n+2}(x)\geq \vphi(n)
\end{matrix}
\right\}
=A_{n,k}\cap I_t.
\]
We conclude
\[
\leb
\left( \bigcup_{k=m+1}^{m+1}\bigcup_{j=M}^{m-1} D_{j,k}\right)
\leq 
\sum_{k=m+1}^{m+1} \leb(A_{n,k}\cap I_t)
\ll 
\frac{\log^{4}\vphi(n)}{\vphi^2(n)}\leb(I_t)
\]
and \eqref{Eq:EstU3} follows.

In view of Lemma \ref{Lem:Divergence:01}, \eqref{Eq:EstU1}, \eqref{Eq:EstU2}, and \eqref{Eq:EstU3}, direct computations give
\begin{align*}
S_1
&\leq 
\sum_{(m,n)\in \clI_1} \leb(U_1(m,n)) +\leb(U_2(m,n))+ \leb(U_3(m,n)) \\
&\ll
\frac{1}{\leb(I_t)} \left(\sum_{n=M}^N \leb(A_n\cap I_t)\right)^2.
\end{align*}
The lemma is now proven.
\end{proof}

\begin{proof}[Proof of Theorem \ref{TEO:MAIN:MEASURE}. Divergence case]
Take any $t\in\Na$ and $\bfb\in\Na^t$. Write $I_t=I_t(\bfb)$ and let $M$ be as above. For each $n\in\Na$ put $E_n:=A_{M-1+n} \cap I_t$. Lemma \ref{Lem:Divergence:01} and the divergence assumption imply $\sum_n \leb(E_n)=\infty$. Then, Lemma \ref{Lem:Divergence:02} and the Chung-Erd{\"o}s inequality yield
\[
\leb\left( \limsup_{n\to\infty} E_n\right)
\gg
\leb(I_t).
\]
Let $\clC$ denote the family of open intervals which are the interior some fundamental interval. The result now follows from $\clF_{3}(\vphi)\cap I_t=\limsup_n E_n$ and Lemma \ref{Le:Knopp}.
\end{proof}


\section{Proof of Theorem \ref{TEO:MAIN:HD}}\label{Section:ProofHD}

The cases $B=1$ and $B=\infty$ follow from $\clF_2(\vphi)\subseteq \clF_3(\vphi)\subseteq \clE_3(\vphi)$ and theorems \ref{TEO:TTW:HD} and \ref{TEO:HWX:HD}. Thus, we assume $1<B<\infty$. Define the functions $X_1,X_2,X_3:[0,1]\to\RE$ by
\[
X_1(s)= \frac{(1-s )^2}{s^2-s+1},  
\quad
X_2(s)= \frac{s^2}{s^2-s+1},  
\quad 
X_3(s)=1-(X_1(s) + X_2(s)).
\]
Direct computations tell us that for all $s\in [0,1]$ we have $0\leq X_1(s),X_2(s),X_3(s)\leq 1$,  and
\begin{align}
g_3(s) &= (2s-1) + X_1(s)s, \label{Eq:EqAux:01} \\
g_3(s) &< (1-s)X_1(s) - s, \label{Eq:EqAux:02} \\
g_3(s) &\leq (3s-1)(X_1(s) + X_2(s)) + s. \label{Eq:EqAux:03}
\end{align}
Write
\[
x_1:=X_1(p_B), \quad
x_2:=X_2(p_B), \quad
x_3:=X_3(p_B).
\]
%In the following, we omit the argument in the partial quotients $a_n$.
\subsection{Upper bound}

We cover $\clF_3(B)$ by three sets determined by the length of the overlap of the restricted blocks of partial quotients.  Define 
\begin{align*}
\clF_{3}^{2}(B)
&:=
\left\{ x\in [0,1): 
\begin{matrix}
a_{n-1}(x) a_{n}(x)a_{n+1}(x) \geq B^n, \\
a_{n}(x)a_{n+1}(x)a_{n+2}(x) \geq B^n 
\end{matrix}
\text{ for i.m.   } n\in\Na \right\}, \\
\clF_{3}^{1}(B)
&:=
\left\{ x\in [0,1): 
\begin{matrix}
a_{n-2}(x) a_{n-1}(x) a_{n}(x) \geq B^n, \\
a_{n}(x) a_{n+1}(x) a_{n+2}(x) \geq B^n 
\end{matrix}
\text{ for i.m.   }n\in\Na\right\}, \\
\clF_{3}^{0}(B)
&:=
\left\{x\in [0,1): 
\begin{matrix}
a_{k}(x)a_{k+1}(x)a_{k+2}(x)\geq B^n, \\
a_{n}(x)a_{n+1}(x)a_{n+2}(x)\geq B^n
\end{matrix}
\text{ for some } 1\leq k\leq n-3 \text{ for i.m.   } n\in\Na \right\}.
\end{align*}
The family $\{\clF_{3}^{2}(B), \clF_{3}^{1}(B), \clF_{3}^{0}(B)\}$ covers $\clF_{3}(B)$. The result will follow from estimating the Hausdorff dimension of each $\clF_{3}^{j}(B)$, $j=0,1,2$. The computations in the three cases are alike. For this reason, we only give a fully detailed proof for $\clF_3^2(B)$ and sketch the argument for $\clF_3^1(B)$ and $\clF_3^0(B)$.
\subsubsection{Overlaps of length 2}
Let us show that 
\begin{equation}\label{Eq:hdim:clF31}
\hdim \clF_{3}^{2}(B) \leq p_B.
\end{equation}
Write $\alpha:=B^{x_1}$. Take $n\in\Na$ and define
\[
\clF_{3,n}^{2}(B)
:=
\left\{ x\in [0,1): 
\begin{matrix}
a_{n-1}(x)a_{n}(x)a_{n+1}(x) \geq B^n, \\
a_{n}(x)a_{n+1}(x)a_{n+2}(x) \geq B^n 
\end{matrix}
\right\}.
\]
We cover $\clF_{3,n}^{2}(B)$ by the sets $\clF_{3,n}^{2,C}(B)$ and $\clF_{3,n}^{2,G}(B)$ given by 
\begin{align*}
\clF_{3,n}^{2,G}(B)
&:=
\left\{ x\in [0,1): a_{n-1}(x) \geq \alpha^n, \, a_{n-1}(x)a_{n}(x) a_{n+1}(x)  \geq B^n, a_{n}(x) a_{n+1}(x) a_{n+2}(x)  \geq B^n \right\}, \\
\clF_{3,n}^{2,C}(B)
&:=
\left\{ x\in [0,1): a_{n-1}(x)< \alpha^n, \, a_{n-1}(x) a_{n}(x) a_{n+1}(x)  \geq B^n \right\}.
\end{align*}
Let us cover $\clF_{3,n}^{2,G}(B)$ with the sets $F'$ and $F''$ defined as 
\begin{align*}
F' &:=
\left\{ x\in [0,1): a_{n-1}(x) \geq \alpha^n, \; \frac{B^n}{a_{n-1}(x)} \leq a_{n}(x) a_{n+1}(x) \leq \frac{B^n}{\alpha^n}, a_{n+2}(x) \geq \frac{B^n}{a_n(x)a_{n+1}(x)} \right\}, \\
F'' &:=
\left\{ x\in [0,1): a_{n-1}(x)\geq \alpha^n, \; \frac{B^n}{\alpha^n} \leq a_{n}(x) a_{n+1}(x) \right\}.
\end{align*}
Put
\begin{align*}
I&:=
\left\{
\bfa=(a_1,\ldots, a_{n+1})\in\Na^{n+1}: a_{n-1}> \alpha^n, \; 1\leq a_n a_{n+1} \leq \frac{B^n}{\alpha^n}\right\}, \\
I' &:=
\left\{
\bfa=(a_1,\ldots, a_{n-1})\in\Na^{n-1}: a_{n-1}> \alpha^n \right\}.
\end{align*}
For any $\bfa\in I$ the set
\[
J_{n+1}(\bfa):
=
\bigcup_{a_{n+2}\geq \frac{B^n}{a_na_{n+1}}} I_{n+2}(\bfa, a_{n+2})
\]
satisfies
\begin{equation}\label{Eq:UpperBound:a}
|J_{n+1}(\bfa)|
\asymp
\sum_{a_{n+2}\geq \frac{B^n}{a_na_{n+1}}} 
\frac{1}{q_{n-1}^2 a_{n}^2a_{n+1}^2 a_{n+2}^2}
\asymp
\frac{1}{q_{n-1}^2 B^n a_{n}a_{n+1}}.
\end{equation}
Similarly, for any $\bfa\in I'$ the set
\[
J_{n-1}'(\bfa)
:=
\bigcup_{a_{n}a_{n+1} \geq \frac{B^n}{\alpha^n } } I_{n+1}(\bfa, a_n, a_{n+1})
\]
satisfies
\begin{equation}\label{Eq:UpperBound:b}
|J_{n-1}'(\bfa)|
\asymp
\frac{1}{q_{n-1}^2}
\sum_{a_{n}a_{n+1} \geq \frac{B^n}{\alpha^n } } 
\frac{1}{a_{n}^2a_{n+1}^2}
\asymp
\frac{\alpha^n}{q_{n-1}^2B^n}n \log\frac{\alpha}{B}
\asymp
\frac{n\alpha^n}{q_{n-1}^2B^n}.
\end{equation}
The last estimate follows from the definition of $\alpha$. Note that 
\[
F'\subseteq\bigcup_{\bfa\in I} J_{n+1}(\bfa)
\quad\text{ and }\quad
F''\subseteq\bigcup_{\bfa\in I'} J_{n-1}'(\bfa).
\]
Now define
\[
I''
:=
\left\{ (a_1,\ldots, a_{n-1})\in \Na^{n-1}: 1\leq a_{n-1}\leq \alpha^n\right\}.
\]
For each $\bfa\in I''$ the set 
\[
J_{n-1}''(\bfa)
:=
\bigcup_{a_{n}a_{n+1}\geq \frac{B^n}{a_{n-1}}} I_{n+1}(\bfa, a_n, a_{n+1})
\]
satisfies
\begin{equation}\label{Eq:UpperBound:c}
|J_{n-1}''(\bfa)|
\asymp
\frac{1}{q_{n-2}^2a_{n-1}B^n}\log \frac{B^n}{a_{n-1}}
\asymp
\frac{1}{q_{n-2}^2a_{n-1}B^n}.
\end{equation}
The last estimate follows from $a_{n-1}\leq \alpha^n$. 

Let $\veps>0$ be arbitrary and write $s=p_B+\veps$. By \eqref{Eq:EqAux:01}, 
\[
(2s-1) + sx_1
> 
(2p_B-1)+x_1p_B
=g_3(p_B)
\]
Since $g_3$ is continuous and strictly increasing, we may pick $\veps_1,\delta_1>0$ satisfying $0<\veps_1<\veps$ and 
\[
(2s-1) + sx_1
>
g_3(p_B+\veps_1) + \delta_1.
\]
Let $N_1\in\Na$ such that $s_n<s + \veps_1<s+\veps$ whenever $n\in\Na_{\geq N_1}$. Then, for such $n$ we have
\[
(2s-1) + sx_1
>
g_3(s_{n-2}) + \delta_1.
\]
As a consequence, taking any $(a_1,\ldots, a_{n-2})\in\Na^{n-2}$ and writing $q_{n-2}=q_{n-2}(a_1,\ldots, a_{n-2})$,
\begin{align*}
\sum_{a_{n-1}\geq \alpha^n} 
\sum_{a_n a_{n+1}\leq \frac{B^n}{\alpha^n}}
 \left( \frac{1}{q_{n-1}^2 B^n a_{n}a_{n+1}} \right)^s
&\asymp 
\frac{1}{q_{n-2}^{2s}B^{ns}}
\sum_{a_{n-1}\geq \alpha^n} 
\frac{1}{a_{n-1}^{2s}}
\sum_{a_n a_{n+1}\leq \frac{B^n}{\alpha^n}}
\frac{1}{a_{n}^{s}a_{n+1}^{s}} \\
&\asymp 
\frac{1}{q_{n-2}^{2s}B^{ns}}
\alpha^{n(1-2s)}
\left(\frac{B}{\alpha}\right)^{n(1-s)} \\
&\asymp 
\frac{1}{q_{n-2}^{2s}}
B^{n(1-2s)} \alpha^{-ns}  \\
&\leq 
\frac{1}{q_{n-2}^{2s_{n-2}} B^{n(g_3(s_{n-2}) + \delta_1)}}.
\end{align*}
This estimate corresponds to \eqref{Eq:UpperBound:a}.

We may use a similar argument relying on \eqref{Eq:EqAux:02} rather than on \eqref{Eq:EqAux:01} to show the existence of some constants $\veps_2,\delta_2>0$ with $0<\veps_2<\veps$ and some $N_2\in\Na$ such that for any $n\in\Na_{\geq N_2}$, any $\bfa=(a_1,\ldots, a_{n-2})\in\Na^{n-2}$ we have%, with $q_{n-2}=q_{n-2}(\bfa)$,
\[
\sum_{a_{n-1}\geq \alpha^n} 
\left(\frac{n\alpha^n}{q_{n-1}^2B^n}\right)^s
\asymp
\frac{n^s}{q_{n-2}^{2s}} \left(\frac{\alpha^n}{B^n}\right)^s\alpha^{n(1-2s)}
\leq
n^s 
\frac{\alpha^{n(1-s)}}{B^{ns}q_{n-2}^{2s}} 
\leq
\frac{1}{q_{n-2}^{2s_{n-2}}B^{n(g_3(s_{n-2}) + \delta_2)}}.
\]
This estimate corresponds to \eqref{Eq:UpperBound:b}.

Likewise, we may show the existence of constants $\veps_3,\delta_3>0$ with $0<\veps_3<\veps$ and $N_3\in\Na$ such that for all $n\in\Na_{\geq N_3}$ and $(a_1,\ldots, a_{n-2})\in\Na^{n-2}$ we have
\[
\sum_{a_{n-1}\leq \alpha^n}
\left(\frac{1}{q_{n-2}^2a_{n-1}B^n}\right)^s
\ll
\frac{1}{q_{n-2}^{2s_{n-2}} B^{n(g_3(s_{n-2}) + \delta_3)}}.
\]
This estimate corresponds to \eqref{Eq:UpperBound:c}.

Therefore, if $N=\max\{N_1,N_2,N_3\}$ and $\delta=\min\{\delta_1, \delta_2, \delta_3\}$, for every $n\in\Na_{\geq N}$ we have
\begin{align*}
\sum_{a_1,\ldots, a_{n-2}\in\Na} 
\sum_{a_n a_{n+1}\leq \frac{B^n}{\alpha^n}}
\left( \frac{1}{B^na_nq_{n-2}^2} \right)^s 
&+ 
\sum_{a_n a_{n+1}\leq \frac{B^n}{\alpha^n}}
 \left( \frac{1}{B^na_nq_{n-2}^2} \right)^s + \\
&\quad+
\sum_{a_{n-1}\leq \alpha^n}
\left(\frac{1}{q_{n-2}^2a_{n-1}B^n}\right)^s 
\ll
B^{-n\delta}.
\end{align*}
We conclude
\[
\clH^{s}(\clF_3^{2}(B))
\leq 
\liminf_{N\to\infty} \sum_{n\geq N} B^{-n\delta} = 0.
\]
Since $\veps>0$ was arbitrary, $\hdim \clF_3^{2}(B)\leq p_B$.
\subsubsection{Overlaps of length 1}
Write $\beta=B^{x_1+x_2}$. Cover  $\clF_{3}^{1}(B)$ by the sets
\begin{align*}
\clF_{3}^{1,C}(B)
&:=
\left\{ x\in \clF_{3}^{1}(B): a_{n-2}(x)a_{n-1}(x) \leq \beta^n \text{ for i.m.  } n\in\Na \right\}\\
\clF_{3}^{1,G}(B,\beta)
&:=
\left\{ x\in \clF_{3}^{1}(B): a_{n-2}(x)a_{n-1}(x)> \beta^n \text{ for i.m.  } n\in\Na \right\}.
\end{align*}
We may cover $\clF_{3}^{1,G}(B)$ by the sets 
\begin{align*}
\clF_{3}^{1,G,1}(B)
&:=
\left\{ x\in[0,1): a_{n-2}(x)a_{n-1}(x)> \beta^n, 1\leq a_n(x)\leq \frac{B^n}{\beta^n}\text{ for i.m.   } n\in\Na \right\}, \\
\clF_{3}^{1,G,2}(B)
&:=
\left\{ x\in[0,1): a_{n-2}(x)a_{n-1}(x)> \beta^n, a_n(x)>\frac{B^n}{\beta^n}\text{ for i.m.   } n\in\Na \right\}.
\end{align*}
For each $n\in \Na$ and every $\bfa=(a_1,\ldots, a_{n-1}, a_n)\in \Na^n$ such that $a_{n-2}a_{n-1}> \beta^n$ and $1\leq a_n\leq \frac{B^n}{\beta^n}$ define 
\[
K_{n}^{G,1}(\bfa)
:=
\bigcup_{a_{n+1}a_{n+2}\geq \frac{B^n}{a_n}} 
I_{n+2}(\bfa, a_{n+1}, a_{n+2}).
\]
Then, using \eqref{Eq:Series:06} and $\beta^n \leq \frac{B^n}{a_n}\leq B^n$, we have 
\[
|K_{n}^{G,1}(\bfa)|
\asymp 
\frac{n}{q_{n-1}^2a_n B^n}.
\]
Considering $a_1,\ldots, a_{n-3}$ fixed, we may use again \eqref{Eq:Series:06} to obtain 
\[
n
\sum_{a_n\leq \frac{B^n}{\beta^n}} 
\sum_{a_{n-2}a_{n-1} \leq \beta^n} 
\frac{1}{q_{n-3}^{2s} a_{n-2}^{2s}a_{n-1}^{2s}a_{n}^{2s} B^{ns} }
\ll
\frac{n}{q_{n-3}^{2s}}B^{n(1-2s)}\beta^{-ns}
\leq 
\frac{n}{q_{n-3}^{2s}}B^{n(1-2s)}\alpha^{-ns}.
\]
As a consequence, arguing as with $\clF_3^2(B)$, we have $\hdim \clF_{3}^{1,G,1}(B) \leq p_B$. In order to estimate $\hdim \clF_{3}^{1,C,2}(B)$, for each $\bfa=(a_1,\ldots, a_{n-1}) \in\Na^{n-1}$ such that $a_{n-2}a_{n-1}\geq \beta^n$ we consider the set
\[
K_{n-1}^{G,2}(\bfa)
:=
\bigcup_{a_{n}\geq \frac{B^n}{\beta^n}} 
I_{n+2}(\bfa, a_{n}),
\]
which satisfies
\[
|K_{n-1}^{G,2}(\bfa)|
\asymp
\frac{\beta^n}{B^nq_{n-1}^2}.
\]
Fixing the terms $a_1,\ldots, a_{n-3}$, for any $s>0$ we have
\[
\sum_{a_{n-2}a_{n-1}\geq \beta^n}
\left( \frac{\beta^n}{B^nq_{n-1}^2}\right)^s 
\asymp
\frac{\beta^{n(1-3s)}}{q_{n-3}^{2s}B^{ns}}.
\]
We may argue as we did for $F''$ using \eqref{Eq:EqAux:03} to conclude $\hdim \clF_{3}^{1,G,2}(B)\leq p_B$.

%The set of those $x=[a_1,a_2,\ldots]\in \clF_3^{1}(B)$ such that for infinitely many $n\in\Na$ check
%\[
%a_{n-2}a_{n-1}<\beta^n, \quad
%a_{n-2}a_{n-1}a_n\geq B^n, 
%\quad\text{ and }\quad
%a_{n}a_{n+1}a_{n+2}\geq B^n
%\]
The set $\clF_3^{1,C}$ is contained in
\[
\clF_{3}^{1,C,1}(B)
:=
\left\{x\in [0,1): a_{n-2}(x) a_{n-1}(x) \leq \beta^n, a_{n}(x) \geq \left(\frac{B}{\beta}\right)^n \text{ for i.m.   }n\in\Na\right\}.
\]
For $n\in\Na$ and $(a_{1},\ldots,a_{n-1})\in\Na^{n-1}$ such that $a_{n-2}a_{n-1}\leq \beta^n$ define 
\[
K_{n-1}(a_1,\ldots, a_{n-1})
=
\bigcup_{a_n> \frac{B^n}{\beta^n}} I_n(a_1,\ldots, a_n),
\]
so
\[
|K_{n-1}(a_1,\ldots, a_{n-1})|
\asymp
\frac{1}{B^nq_{n-3}^s a_{n-1}a_{n-2}}.
\]
Moreover, fixing $a_1,\ldots, a_{n-3}$ and assuming that $n$ is large, every $s>0$ verifies
\[
\sum_{a_{n-1}a_{n-2}< \frac{B^n}{\beta^n}}
\left(\frac{1}{B^nq_{n-3}^s a_{n-1}a_{n-2}}\right)^s
\asymp
\frac{B^{n(1-2s)}\beta^{-ns}}{q_{n-3}^{2s}}
<
\frac{B^{n(1-2s)}\alpha^{-ns}}{q_{n-3}^{2s}}.
\]
Arguing as above, we conclude that $\hdim \clF_{3}^{1,C,1}(B)\leq p_B$ and, thus, $\hdim \clF_{3}^{1,C}(B)\leq p_B$ .


\subsubsection{No overlaps}
We content ourselves with describing the useful covering of $\clF_3^{0}(B)$. First, for $n\in\Na$, $k\in \{1,\ldots, n-3\}$, and $\bfa\in\Na^{n-1}$ such that $a_{k}a_{k+1}a_{k+2}\geq B^n$ define 
\[
K_{n-1}^{3,G,k}(\bfa)
:=
\bigcup_{a_n\geq \alpha^n} I_n(\bfa, a_n), 
\quad\text{ so }\quad
|K_{n-1}^{3,G,k}(\bfa)|
\asymp 
\frac{1}{\alpha^n q_{n-1}^2}.
\]
Therefore, letting $a_k,a_{k+1},a_{k+2}$ run along the pairs satisfying $a_ka_{k+1}a_{k+2}>B^n$ considering the remaining terms fixed, we have
\[
\sum_{a_ka_{k+1}a_{k+2}>B^n} 
\frac{1}{\alpha^{ns} q_{n-1}^{2s}}
\ll
\frac{B^{n(1-2s)}}{\alpha^{ns} q_{n-4}^{2s}}
\]
and we proceed as in the estimate for $\clF_3^2(B)$.


Second, for any $n\in\Na$, $k\in \{1,\ldots, n-3\}$, $\bfa\in\Na^n$ such that $a_{k}a_{k+1}a_{k+2}\geq B^n$ and $a_n\leq \alpha^n$ define 
\[
K_n^{3,C,k}(\bfa)
:=
\bigcup_{a_{n+1}a_{n+2}\geq \frac{B^n}{a_n}} I_n(\bfa, a_{n+1}, a_{n+ 2}), 
\quad\text{ so }\quad
|K_n^{3,C,k}(\bfa)|
\asymp 
\frac{n}{B^na_nq_{n-1}^2}.
\]
Hence, varying $a_k,a_{k+1}$ along the pairs satisfying $a_ka_{k+1}\geq B^n$ and $a_n$ along $1\leq a_n\leq \alpha^n$, for all $s>0$ we have 
\[
\sum_{a_ka_{k+1}\geq B^n}
\sum_{a_n\leq \alpha^n}
\left(\frac{n}{B^na_nq_{n-1}^2}\right)^s
\ll
\frac{n^s}{B^{ns}q_{n-3}^{2s}}B^{n(1-2s)}\alpha^{n(1-s)}
\leq
\frac{n^s}{q_{n-3}^{2s}}B^{n(1-3s)}\alpha^{n(1-s)}.
\]
Direct computations tell us that for $s\geq \frac{1}{2}$ we have
\[
(3s-1)+X_1(s)(s-1)>g_3(s).
\]
Arguing as we did for $\clF_3^2(B)$, we may conclude that $\hdim \clF_3^0(B)\leq p_B$.

\subsection{Lower bound}

Let $x_1,x_2,x_3$ be as in the upper bound and define
\[
A_0 = B^{x_1}, \quad
A_1 = B^{x_2}, \quad
A_2 = B^{x_3}, \quad
A_3 = B^{x_1}.
\]
With notation of Lemma \ref{Le:MumtazNikita}, we have $S_4(A_0,A_1,A_2,A_3)\subseteq \clF_3(B)$ and
\[
\hdim S_4(A_0,A_1,A_2,A_3)\leq  \hdim \clF_3(B).
\]
Also, $A_0A_1A_2=A_1A_2A_3=B$ and
\[
\log \beta_{-1} =0, \quad
\log \beta_{0} = x_1\log B, \quad
\log \beta_{1} = (x_1+x_2)\log B, 
\]
\[
\log \beta_{2} = \log B, \quad
\log \beta_{3} = (1+x_1)\log B.
\]
As a consequence, the potentials $\vphi_{i,s}$ defining $d_i$ are
\begin{align*}
\vphi_{0,s}(x) 
&= 
-sx_1\log B -s \log|T'(x)|, \\
\vphi_{1,s}(x) 
&= 
-((2x_1+x_2)s - x_1) \log B - s\log|T'(x)|, \\
\vphi_{2,s}(x) 
&= 
-((1+x_1+x_2)s-x_1-x_2) \log B - s\log|T'(x)|, \\
\vphi_{3,s}(x) 
&= -((2+x_1)s-1)\log B  - s\log|T'(x)|.
\end{align*}
It is not hard to show that $p_B=d_1=d_2=d_3\leq d_0$. Certainly, since $sX_1(s)\leq (2+X_1(s))s-1$ when $\frac{1}{2}\leq s\leq 1$, we have $d_3\leq d_0$. We only explain $p_B=d_1$, the other equalities are handled in a similar fashion. First, note that 
\[
g_3(p_B)
= 
(2+x_1)p_B-1
\]
and that the functions $g_3$ and $s\mapsto (2+x_1)s-1$ are strictly increasing and continuous on $\left[\frac{1}{2}, 1\right]$. If $\veps>0$ and $s=p_B+\veps$, then 
\[
(2+x_1)s-1
>
g(p_B)
\]
and we would have $P(T;\vphi_{1,s})\leq 0$, hence $d_1\leq p_B+\veps$. Since $\veps$ was arbitrary, $d_1\leq p_B$. We now show that $d_1< p_B$ leads to a contradiction. If the strict inequality were true, we would be able to pick some number $s$ with $d_1<s<p_B$ and close enough to $p_B$ such that $g(s)=(2+x_1)s-1<g(p_B)$. However, for such $s$ we would have 
\[
0>
P(T;\vphi_{1,s})
= 
P(T;-g_3(s)\log B - s\log|T'|)>0.
\]
By virtue of Lemma  \ref{Le:MumtazNikita}, $p_B\leq \hdim \clF_3(B)$.

\section{Final remarks and further results}\label{Section:FinalRemarks}

We stated in the introduction that Theorem \ref{TEO:MAIN:MEASURE} is an important piece in the extension of the laws of large numbers due to Diamond and Vaaler (Equation \eqref{Eq:Thm:DV}) and Hu, Hussain, and Yu (Equation \eqref{Eq:Thm:HHY}). In a forthcoming article, we study the asymptotic behaviour of the sums of products of consecutive partial quotients. More precisely, for $n,\ell\in\Na$ and any irrational number $x\in [0,1)$ define
\[
S_{n,\ell}(x)=\sum_{j=1}^n a_{j}(x) \cdots a_{j+\ell-1}(x).
\]
\begin{teo01}\label{TEO:PROMISED:01}\
For every $\veps>0$ we have
\[
\lim_{n\to \infty}
\leb
\left(\left\{ 
x\in [0,1): 
\left| \frac{1}{n\log^{\ell}n} S_{n,\ell}(x) - \frac{1}{\ell \log 2}\right|\geq \veps\right\}\right)=0.
\]
Almost every $x\in [0,1)$ satisfies
\[
\lim_{n\to\infty} 
\frac{1}{n\log^{\ell} n} 
\left( S_{n,\ell}(x) - \max_{1\leq j \leq n} a_{j}(x)\cdots a_{j+\ell-1}(x)\right)
=
\frac{1}{\ell \log 2}.
\]
\end{teo01}

These results could be extended to products of partial quotients in arithmetic progressions. Given any $d\in\Na$,  for $n\in\Na$  and any irrational $x\in [0,1)$ define the sum 
\[
S_{n,\ell}^d(x)
=
\sum_{j=1}^n a_{j}(x)a_{j+d}(x) \cdots a_{j+(\ell-1)d}(x).
\]
\begin{propo01}
Theorem \ref{TEO:PROMISED:01} holds if $S_{n,\ell}(x)$ is replaced with $S_{n,\ell}^d(x)$.
\end{propo01}


Regarding the Hausdorff dimension of certain Lebesgue null sets, for an increasing function $\phi:\Na\to\RE$ with $\phi(n)\to\infty$ as $n\to\infty$ we consider the set 
\[
E_{\ell}(\phi)
:=
\left\{ x\in [0,1): \lim_{n\to\infty} \frac{S_{n,\ell}(x)}{\phi(n)}=1\right\}.
\]
\begin{teo01}\ 
Let $\alpha$ be a positive real number.
\begin{enumerate}[\rm (I)]
    \item If $0<\alpha<\frac{1}{2}$ and $\phi(n)=e^{n^{\alpha}}$, then $\hdim E_{\ell}(\phi)=1$.
    \item If $\frac{1}{2}\leq \alpha$ and $\phi(n)=e^{n^{\alpha}}$, then $\hdim E_{\ell}(\phi)=\frac{1}{2}$.
    \item If $1<\alpha$ and $\phi(n)=e^{\alpha^n}$, then $\hdim E_{\ell}(\phi)=\frac{1}{1+\alpha}$.
\end{enumerate}
\end{teo01}
 Recently, Song, Yu, and Yu in \cite{SYY2025} investigated the Hausdorff dimension of a related set concerning the product of two partial quotients. 
 More precisely, for any $x\in [0,1)\setminus\QU$ and $n\in\Na$ write.
 \[
 L_n(x) = \max_{1\leq i\leq n} \, a_i(x)a_{i+1}(x),
 \]
  Given a non-decreasing function $\vphi:\RE_{>0}\to\RE_{>0}$ satisfying certain regularity properties, they computed the Hausdorff dimension of 
 %they investigated the set for $L_n(x) = \max_{1\leq i\leq n}\{a_i(x)a_{i+1}(x)\}$:
        \[
            L(\varphi):= \left\{ x\in(0,1): \lim_{n\to \infty} \frac{L_n(x)}{\varphi(n)}=1\right\}.
        \] 
        Their work lead us to the following problem.
        \begin{prob01}
        Let $\ell\in\Na_{\geq 3}$.
        Can we extend the results in \cite{SYY2025} to the product of $\ell\in\Na_{\geq 3}$ consecutive partial quotients? 
        That is, for $x\in [0,1)\setminus\QU$ and $n\in\Na$, define
        \[
        L_{\ell,n}(x) = \max_{1\leq i\leq n}\, a_i(x)a_{i+1}(x)\cdots a_{i+ \ell -1 }(x) .
        \]
        Given $\rho>0$ and a function $\vphi:\RE_{>0}\to\RE_{>0}$ such that $\log \vphi$ is regularly increasing of index $\rho$, determine the Hausdorff dimension of the set
        \[
            L_{\ell}(\varphi)
            := 
            \left\{ x\in(0,1): \lim_{n\to \infty} \frac{L_{\ell,n}(x)}{\varphi(n)}=1\right\}.
        \] 
        \end{prob01}
        %It remains to extend their paper to the product of a block arbitrary length of consecutive partial quotients rather than the block of just two.         
        %There problem with regard to investigating the set involving an arbitrary length of consecutive partial quotients rather than the block of two is still open.
\begin{rema01}
The $n$-th term of one of the series in Theorem \ref{TEO:MAIN:MEASURE} is
\[
\frac{n\log^{4}\vphi(n)}{\vphi^2(n)}
+
\frac{\log \vphi(n)}{\vphi(n)}.
\]
The first term accounts for the large products appearing in non-overlapping blocks and the second term corresponds overlapping blocks. Hence, for $\ell=1$ the second term is non-existent. %Nevertheless, it is an interesting problem by itself to check whether Theorem \ref{TEO:MAIN:MEASURE} still holds for $\ell=1$. This amounts to prove that the convergence of the series $\sum_n\frac{n}{\vphi^2(n)}$ implies that of $\sum_n (\vphi \log \vphi(n))^{-1} $. 
\end{rema01}

\begin{rema01}
We suspect that arguments resembling the proof of Theorem \ref{TEO:MAIN:HD} can lead to the computation of $\hdim \clF_{\ell}(\vphi)$ for $\ell\geq 4$. This could provide a recursive definition of the potential giving the Hausdorff dimension when $1<B<\infty$ (see Theorem \ref{TEO:HWX:HD}). However, we stress that the proofs for $\ell\in\{2,3\}$ do not rely on the result for lower values of $\ell$.
\end{rema01}





\bibliography{References}
\bibliographystyle{abbrv}
\end{document}



\begin{thebibliography}{10}

\bibitem{BHS}
P.~Bos, M.~Hussain, and D.~Simmons.
\newblock The generalised {H}ausdorff measure of sets of {D}irichlet
  non-improvable numbers.
\newblock {\em Proc. Amer. Math. Soc.}, 151(5):1823--1838, 2023.

\bibitem{BugeaudBookCUP160}
Y.~Bugeaud.
\newblock {\em Approximation by algebraic numbers}, volume 160 of {\em
  Cambridge Tracts in Mathematics}.
\newblock Cambridge University Press, Cambridge, 2004.

\bibitem{ChungErdos1952}
K.~L. Chung and P.~Erd\"{o}s.
\newblock On the application of the {B}orel-{C}antelli lemma.
\newblock {\em Trans. Amer. Math. Soc.}, 72:179--186, 1952.

\bibitem{DiamondVaaler1986}
H.~G. Diamond and J.~D. Vaaler.
\newblock Estimates for partial sums of continued fraction partial quotients.
\newblock {\em Pacific J. Math.}, 122(1):73--82, 1986.

\bibitem{EinsiedlerWardET}
M.~Einsiedler and T.~Ward.
\newblock {\em Ergodic theory with a view towards number theory}, volume 259 of
  {\em Graduate Texts in Mathematics}.
\newblock Springer-Verlag London, Ltd., London, 2011.

\bibitem{FalconerBookFractal}
K.~Falconer.
\newblock {\em Fractal geometry}.
\newblock John Wiley \& Sons, Ltd., Chichester, third edition, 2014.
\newblock Mathematical foundations and applications.

\bibitem{FengWuLiang1997}
D.~J. Feng, J.~Wu, J.-C. Liang, and S.~Tseng.
\newblock Appendix to the paper by {T}. \l uczak---a simple proof of the lower
  bound: ``{O}n the fractional dimension of sets of continued fractions''.
\newblock {\em Mathematika}, 44(1):54--55, 1997.

\bibitem{HanusMauldinUrbanski2002}
P.~Hanus, R.~D. Mauldin, and M.~Urba\'{n}ski.
\newblock Thermodynamic formalism and multifractal analysis of conformal
  infinite iterated function systems.
\newblock {\em Acta Math. Hungar.}, 96(1-2):27--98, 2002.

\bibitem{HardyLittlewood1923}
G.~H. Hardy and J.~E. Littlewood.
\newblock The {A}pproximate {F}unctional {E}quation in the {T}heory of the
  {Z}eta-{F}unction, with {A}pplications to the {D}ivisor-{P}roblems of
  {D}irichlet and {P}iltz.
\newblock {\em Proc. London Math. Soc. (2)}, 21:39--74, 1923.

\bibitem{HensleyBook}
D.~Hensley.
\newblock {\em Continued fractions}.
\newblock World Scientific Publishing Co. Pte. Ltd., Hackensack, NJ, 2006.

\bibitem{HuHussainYu2021}
H.~Hu, M.~Hussain, and Y.~Yu.
\newblock Limit theorems for sums of products of consecutive partial quotients
  of continued fractions.
\newblock {\em Nonlinearity}, 34(12):8143--8173, 2021.

\bibitem{HuangWuXu2020}
L.~Huang, J.~Wu, and J.~Xu.
\newblock Metric properties of the product of consecutive partial quotients in
  continued fractions.
\newblock {\em Israel J. Math.}, 238(2):901--943, 2020.

\bibitem{HKWW}
M.~Hussain, D.~Kleinbock, N.~Wadleigh, and B.-W. Wang.
\newblock Hausdorff measure of sets of {D}irichlet non-improvable numbers.
\newblock {\em Mathematika}, 64(2):502--518, 2018.

\bibitem{ShulgaHussain2023Exp}
M.~Hussain and N.~Shulga.
\newblock Metrical properties of exponentially growing partial quotients.
\newblock {\em Forum Mathematicum}, 2024.

\bibitem{ShulgaHussainPisa}
M.~Hussain and N.~Shulga.
\newblock Metrical properties of finite product of partial quotients in
  arithmetic progressions.
\newblock {\em In press: Ann. Sc. Norm. Super. Pisa Cl. Sci.}, 2024.

\bibitem{KallenbergBook2021}
O.~Kallenberg.
\newblock {\em Foundations of modern probability}, volume~99 of {\em
  Probability Theory and Stochastic Modelling}.
\newblock Springer, Cham, third edition, [2021] \copyright 2021.

\bibitem{KhinchinCFBook}
A.~Y. Khinchin.
\newblock {\em Continued fractions}.
\newblock University of Chicago Press, Chicago, Ill.-London, 1964.

\bibitem{Khintchine1935}
A.~Khintchine.
\newblock Metrische {K}ettenbruchprobleme.
\newblock {\em Compositio Math.}, 1:361--382, 1935.

\bibitem{KleinbockWadleigh2019}
D.~Kleinbock and N.~Wadleigh.
\newblock An inhomogeneous {D}irichlet theorem via shrinking targets.
\newblock {\em Compos. Math.}, 155(7):1402--1423, 2019.

\bibitem{Knopp1926}
K.~Knopp.
\newblock Mengentheoretische behandlung einiger probleme der diophantischen
  approximationen und der transfiniten wahrscheinlichkeiten.
\newblock {\em Mathematische Annalen}, 95:409--426, 1926.

\bibitem{LiWangWuXu2014}
B.~Li, B.-W. Wang, J.~Wu, and J.~Xu.
\newblock The shrinking target problem in the dynamical system of continued
  fractions.
\newblock {\em Proc. Lond. Math. Soc. (3)}, 108(1):159--186, 2014.

\bibitem{Luc1997}
T.~{\L}uczak.
\newblock On the fractional dimension of sets of continued fractions.
\newblock {\em Mathematika}, 44(1):50--53, 1997.

\bibitem{MauldinUrbanskiBook}
R.~D. Mauldin and M.~Urba\'{n}ski.
\newblock {\em Graph directed {M}arkov systems}, volume 148 of {\em Cambridge
  Tracts in Mathematics}.
\newblock Cambridge University Press, Cambridge, 2003.
\newblock Geometry and dynamics of limit sets.

\bibitem{Philipp1988}
W.~Philipp.
\newblock Limit theorems for sums of partial quotients of continued fractions.
\newblock {\em Monatsh. Math.}, 105(3):195--206, 1988.

\bibitem{SYY2025}
K.~Song, D.~Yu, and Y.~Yu.
\newblock Multifractal analysis of maximal product of consecutive partial
  quotients in continued fractions.
\newblock {\em arXiv:2502.02064}, 2025.

\bibitem{Sprindzuk1979}
V.~G. Sprind\v{z}uk.
\newblock {\em Metric theory of {D}iophantine approximations}.
\newblock Scripta Series in Mathematics. V. H. Winston \& Sons, Washington, DC;
  John Wiley \& Sons, New York-Toronto-London, 1979.
\newblock With a foreword by Donald J. Newman, A Halsted Press Book.

\bibitem{TanTianWang2023}
B.~Tan, C.~Tian, and B.~Wang.
\newblock The distribution of the large partial quotients in continued fraction
  expansions.
\newblock {\em Sci. China Math.}, 66(5):935--956, 2023.

\bibitem{TanZhou2024}
B.~Tan and Q.-L. Zhou.
\newblock Metrical properties of the large products of partial quotients in
  continued fractions.
\newblock {\em Nonlinearity}, 37(2):Paper No. 025008, 28, 2024.

\bibitem{WangWu2008}
B.-W. Wang and J.~Wu.
\newblock Hausdorff dimension of certain sets arising in continued fraction
  expansions.
\newblock {\em Adv. Math.}, 218(5):1319--1339, 2008.

\bibitem{Wu2006}
J.~Wu.
\newblock A remark on the growth of the denominators of convergents.
\newblock {\em Monatsh. Math.}, 147(3):259--264, 2006.

\end{thebibliography}
